\documentclass[11pt]{amsart}
\IfFileExists{lmodern.sty}%
  {\usepackage[T1]{fontenc}\usepackage{lmodern}\usepackage{microtype}}%
  {\usepackage[expansion=false]{microtype}}
\usepackage{amssymb,mathtools}
\usepackage{booktabs,array}
\usepackage{enumitem}
\emergencystretch=2.5em
\usepackage[colorlinks=true,linkcolor=blue,citecolor=blue,urlcolor=blue]{hyperref}

\newcommand{\FormalizationURL}{https://github.com/slitsey/catalan-irrationality-measure}
\newcommand{\FormalizationLink}{\url{\FormalizationURL}}

% ---------------------------------------------------------------------------
% Numerical constants. Each is the decimal truncation (or outward rounding) of a rational
% number of the Lean proof; the Lean names are in the comments.
\newcommand{\Bx}{409.99999}                             % 53 rho'/(eta - 212 delta)
\newcommand{\Rho}{296.586094}                           % rho' = K590.rhoEnd (879663/10000)
\newcommand{\RhoOld}{301.54258}                         % 58 c_2 + w (used only in the remark on rho)
\newcommand{\Ctwo}{3.596893}                            % secondCost'
\newcommand{\Cone}{3.549542}                            % firstCost'
\newcommand{\PsiZero}{87.9663}                          % Psi_0, weight plus binomial per column
\newcommand{\Utarget}{2.255858836}                      % K590.target
\newcommand{\Gain}{8.29824}                             % 58^2 (target0 - target)
\newcommand{\Uzero}{2.2583256}                          % target0
\newcommand{\NormCap}{0.7736205195}                     % normCap
\newcommand{\SupX}{-0.96450499}
\newcommand{\SupY}{-1.57499400}
\newcommand{\FieldCap}{-2.53949899}                     % fieldCap
\newcommand{\ResCap}{-0.0631062}                        % residualCap
\newcommand{\EtaInf}{7.98775}                           % real place (old) minus denominators, no gain
\newcommand{\EtaTot}{38.33924}                          % eta of the criterion
\newcommand{\Kappa}{1.382397}                           % 410/rho'
\newcommand{\Kcoef}{160}
\newcommand{\Ngeo}{75}
\newcommand{\Ncomplex}{41}
\newcommand{\Nblocks}{47, 38 and 23}
\newcommand{\Ncells}{373, 302 and 176}
\newcommand{\Mmargin}{2.986\cdot10^{-6}}
\newcommand{\TDelta}{6431.1338}                         % 58^2 (target0 - log 2 / 2)
\newcommand{\TDeltaOld}{6422.8355}                      % 58^2 (target - log 2 / 2)
\newcommand{\EtaExact}{38.34580}                        % TDelta minus the denominator exponent
\newcommand{\PayL}{53.0000849}                          % d + 4 d delta kappa
\newcommand{\PayR}{53.0000857}                          % eta kappa
\newcommand{\Digest}{ff307ca0287dc5224e268202cb0b09776e933652acd2f2f879492af28d3e381d}
\newcommand{\Chi}{22.05972}                             % area of the rank-deficiency profile
\newcommand{\ChiLean}{22.05325}                         % gain532_ge
\newcommand{\Aden}{6392.7880}                           % exponent of the denominators
\newcommand{\Wcap}{52.788}
\newcommand{\NcellsSigned}{\textbf{??}}                 % filled in after the cell build
% ---------------------------------------------------------------------------

\newtheorem{theorem}{Theorem}[section]
\newtheorem{lemma}[theorem]{Lemma}
\newtheorem{proposition}[theorem]{Proposition}
\newtheorem{corollary}[theorem]{Corollary}
\theoremstyle{definition}
\newtheorem{definition}[theorem]{Definition}
\theoremstyle{remark}
\newtheorem{remark}[theorem]{Remark}

\numberwithin{equation}{section}
\setcounter{tocdepth}{1}

\newcommand{\Z}{\mathbb Z}
\newcommand{\Q}{\mathbb Q}
\newcommand{\R}{\mathbb R}
\newcommand{\N}{\mathbb N}
\newcommand{\C}{\mathbb C}
\newcommand{\F}{\mathbb F}
\newcommand{\rank}{\operatorname{rank}}
\newcommand{\cV}{\mathcal V}
\newcommand{\cR}{\mathcal R}
\newcommand{\cA}{\mathcal A}
\newcommand{\cL}{\mathcal L}
\newcommand{\cK}{\mathcal K}
\newcommand{\rev}{\operatorname{rev}}
\newcommand{\lean}[1]{\texttt{#1}}

\begin{document}

\title{An irrationality measure for Catalan's constant}

\author{Sean Litsey}
\email{sean.litsey.math@gmail.com}

\date{November 1, 2026}

\subjclass[2020]{Primary 11J82; Secondary 11J72, 15A15, 31A15, 68V20}

\keywords{Catalan's constant, irrationality measure, determinants, rank of reductions modulo p,
logarithmic potentials, Pick matrices, Heine's formula, Gauss quadrature, formal verification}

\begin{abstract}
Let $G=\sum_{j\ge0}(-1)^j(2j+1)^{-2}$ be Catalan's constant. We prove that $|G-a/b|>b^{-410}$ for
every integer $a$ and every sufficiently large integer $b$. Thus the irrationality exponent of $G$ is
at most $410$. The proof builds on OpenAI's proof that $G$ is irrational. That proof compares two
bounds for a determinant whose entries are linear forms in $1$ and $G$. The lower bound comes from
the denominators, and the upper bound from an integral representation. We improve both sides and add
the uniformity a measure needs. At the finite places we bound the rank of the reduction modulo $p$ of
the matrix of rational parts. This lowers the denominator at every prime in an explicit window. At the
real place we bound the determinant of the row functions by a signed estimate. A Pick matrix argument
reduces it to the block of the negative nodes. Heine's formula writes that block as a ratio of two
log-gas integrals, and both are bounded by logarithmic energy. A Gauss quadrature formula extends the
bound to an exponentially small neighbourhood of $G$, at a cost paid column by column. All steps are
formalized in the Lean proof assistant, with only the standard axioms.
\end{abstract}

\maketitle

\tableofcontents

\section{Introduction}\label{sec:intro}

\subsection{The result}
Catalan's constant
\[
 G=\sum_{j=0}^\infty\frac{(-1)^j}{(2j+1)^2}=L(2,\chi_{-4})=0.91596\,55941\,77219\ldots
\]
is the value at $s=2$ of the Dirichlet $L$-function of the nontrivial character modulo $4$. Its
arithmetic has long been studied by analogy with $\zeta(2)$ and $\zeta(3)$. Until recently it was not
known to be irrational. In October 2026 OpenAI \cite{OpenAI2026} proved that $G$ is irrational and
verified the proof in Lean.

The \emph{irrationality exponent} $\mu(\xi)$ of a real number $\xi$ is the infimum of the real
numbers $\mu$ such that $|\xi-a/b|>b^{-\mu}$ for all integers $a$ and all sufficiently large
integers $b$. Every irrational number has $\mu(\xi)\ge2$, and almost every real number has
$\mu(\xi)=2$. An \emph{irrationality measure} is an explicit upper bound for $\mu(\xi)$.

\begin{theorem}\label{thm:main}
For every integer $a$ and every sufficiently large integer $b$,
\[
 \left|G-\frac ab\right|>b^{-410}.
\]
In particular $\mu(G)\le410$.
\end{theorem}

The proof uses the prime number theorem. The theorem and every result it depends on have been verified
in the Lean~4 proof assistant on top of OpenAI's formalization. The formal proof has no unproved steps
and uses only the standard axioms (Section~\ref{sec:formal}). The use of AI in this work is disclosed at
the end of the paper.

\subsection{Background}
Irrationality proofs for values of $L$-functions at integers almost always follow the pattern of
Ap\'ery's proof for $\zeta(3)$ \cite{Apery79}. One constructs integer linear forms $q_n\xi-p_n$ that
tend to zero, with $q_n\ne0$, faster than the growth of their coefficients would allow for a rational
$\xi$. If $|q_n\xi-p_n|=e^{-\alpha n+o(n)}$ and $|q_n|=e^{\beta n+o(n)}$, then $\xi$ is irrational as
soon as $\alpha>0$, and $\mu(\xi)\le1+\beta/\alpha$. Beukers \cite{Beukers79} gave the linear forms for
$\zeta(2)$ and $\zeta(3)$ as integrals of Legendre-type polynomials. Hata \cite{Hata93} refined such
integrals. He used systematically the primes that divide all the coefficients of the forms, an idea that
goes back to Chudnovsky. Rhin and Viola \cite{RhinViola96,RhinViola01} found that families of
integrals for $\zeta(2)$ and $\zeta(3)$ are permuted by a finite group. The invariance of the integrals
lets one divide all their coefficients by a large common factor. This \emph{arithmetic method} gave
the sharpest measures for $\zeta(2)$ and $\zeta(3)$ for many years. Marcovecchio \cite{Marcovecchio09}
applied it to $\log2$. Zeilberger and Zudilin \cite{ZeilbergerZudilin20} obtained
$\mu(\pi)\le7.1032\ldots$ from Salikhov's integrals \cite{Salikhov08} with an additional arithmetic
saving.

For Catalan's constant the analogous linear forms exist but do not decay fast enough. Rivoal and Zudilin
\cite{RivoalZudilin2003} proved that one of the numbers $\beta(2),\beta(4),\ldots,\beta(2m)$ is
irrational for suitable $m$, where $\beta$ is the Dirichlet beta function. This does not single out
$G=\beta(2)$. Zudilin \cite{Zudilin2003} found an Ap\'ery-like recurrence for $G$. Its solutions give
linear forms that tend to zero, but too slowly to prove irrationality. Related constructions and their
limits are described in \cite{Zudilin02Cheb,KrattRivoal08,KrattZudilin19,Zudilin19Survey}. Eskandari,
Murty and Nemoto \cite{EMN2025} construct a mixed motive with period $G$ and obtain linear forms in $1$
and $G$ from it. Eskandari \cite{Eskandari2026} improves the approximation exponent of these forms. Leibtag
et al.\ \cite{LBKWSK26} found further Ap\'ery-like approximations to $G$ in conservative matrix fields.
Numerically the best of them still has a negative approximation exponent.

OpenAI's argument \cite{OpenAI2026} has a different shape. It uses a determinant $\Delta$ of size $48N$
instead of a sequence of linear forms. The entries of $\Delta$ are rational linear combinations of $1$
and $G$, given by double integrals against the Cauchy kernel $1/(1-ts)$. If $G$ were rational, the
denominators of the entries would give a lower bound $|\Delta|\ge e^{-(\mathsf A+o(1))n^2}$ with $n=48N$.
An integral representation gives the upper bound $|\Delta|\le e^{-(\mathsf U+o(1))n^2}$. It rests on
Andr\'eief's identity \cite{Forrester2018Andreief}, familiar from random matrix theory, and an energy
estimate shows $\mathsf U>\mathsf A$. The margin is small, about $1.3\cdot10^{-4}$. The nonvanishing of
$\Delta$ along the primes $N=p$ is proved separately by reduction modulo $p$. Determinants of linear forms
were used for irrationality before, in Zudilin's determinantal criterion with Hankel determinants
\cite{Zudilin}.

The first irrationality measure for $G$ was obtained by Jha \cite{Jha2026}, who proved
$\mu(G)\le10^8$. Carreira \cite{Carreira2026} treated the five integer parameters of the construction
as free. He stated conditions under which OpenAI's argument goes through for general parameters and
obtained $\mu(G)<2294$. He also formulated the best real-place constant of the method as an equilibrium
problem for two measures. In \cite{Carreira2026b} he records attempts to carry the method to $\beta(4)$.
Bojnal \cite{Bojnal2026} instantiated Carreira's general theorem at the parameters $(36,6,2,3,2)$ and
obtained $\mu(G)<2401.74$. The search and certification were carried out with Claude Fable~5.1. Xu
\cite{Xu2026} reported $\mu(G)\le727.03$ for the parameters $(20,3,1,2,1)$, with an improved trial
potential in the energy estimate. His certificate is numerical, in ball arithmetic, and is not
formalized. Sun's preprint \cite{Sun2026} had claimed a proof of the irrationality of $G$. Wachs
\cite{Wachs2026} reports that its main estimate is not justified.

Since September 2026 a series of new results in this area has appeared, several of them obtained with AI
systems. OpenAI proved $\mu(\pi)=2$, with a Lean formalization \cite{OpenAIPi2026}. Liu, Jiang and
Zhang proved $\mu(\log r)=2$ for every rational $r>0$ with $r\ne1$ \cite{LiuJiangZhang26}. Giraudin
lowered the bound for $\mu(\pi^2)$ to $5.0193784$ with two hypergeometric constructions of Zudilin
\cite{Zudilin14} at new parameters \cite{Giraudin26}. He also lowered the bound for $\mu(\zeta(3))$ to
$5.5138800$ \cite{Giraudin26z3}. Leibtag, Barkan, Kalisch, Weinbaum, Shalyt and Kaminer \cite{LBKWSK26}
showed that the classical constructions are one-dimensional trajectories in conservative matrix fields.
Searching over other trajectories, they obtained $\mu(\pi^2)<4.62675$ and $\mu(\zeta(3))<5.4197$.
Prellberg \cite{Prellberg26} proved that $\mathrm{Li}_2(-1/4)$, $\mathrm{Li}_2(1/5)$ and
$\mathrm{Li}_2(-1/3)$ are irrational. Fauzan announced a proof that $\zeta(5)$ is irrational, using
Hankel determinants \cite{Fauzan2026}. Lai combined this method with simultaneous approximations to show
that $1$, $\zeta(3)$ and $\zeta(4)$ are linearly independent over $\Q$ \cite{Lai2026}. For
$L(2,\chi_{-3})$, Calegari, Dimitrov and Tang proved that $1$, $\pi^2$ and $L(2,\chi_{-3})$ are linearly
independent over $\Q$ \cite{CDT2024}. Their effective holonomy bounds give $\mu(L(2,\chi_{-3}))\le24781$
\cite[Thm.~5.1]{CDT2025}.

\subsection{The companion paper}
Catalan's constant is a difference of two values of the trigamma function $\psi'$. For $0<s<t\le1$ put
\[
 \Psi(s,t)=\psi'(s)-\psi'(t)=\sum_{n=0}^\infty\Bigl(\frac1{(n+s)^2}-\frac1{(n+t)^2}\Bigr).
\]
Then $G=\frac1{16}\Psi(\frac14,\frac34)$, and also $L(2,\chi_{-3})=\frac19\Psi(\frac13,\frac23)$ and
$\pi^2=3\Psi(\frac12,1)$. The irrationality exponent does not change under $\xi\mapsto r\xi+r'$ with
$r,r'\in\Q$ and $r\ne0$. So a measure for $G$ is a measure for $\Psi(\frac14,\frac34)$. The companion
paper \cite{Litsey2026b} proves irrationality measures for the values $\Psi(s,1)$ with $s\in\{\frac14,
\frac34,\frac13,\frac23,\frac12\}$ and $s=\frac k5$, $\frac k7$, and for $\Psi(\frac13,\frac23)$. Among
them are $\mu(G\pm5\pi^2/48)$ and $\mu(L(2,\chi_{-3}))$. Its determinants are L\"owner matrices of digamma
differences. That method covers the pair $(\frac14,\frac34)$ as well, but its bound for this pair is too
weak to give a measure. The present paper therefore refines OpenAI's construction instead. The two papers
share their notation, their literature and the formalization repository.

\subsection{What a measure requires}\label{sec:intro-need}
Write $\Delta_N(X)$ for the determinant regarded as a polynomial in the constant, so that $\Delta_N(G)$
is OpenAI's determinant. Let $Q_N$ be the integer polynomial obtained from $\Delta_N$ by clearing
denominators. Three properties of $Q_N$ turn OpenAI's comparison into a measure.
\begin{enumerate}[label=(\roman*)]
\item \emph{Smallness on a neighbourhood of $G$.} The upper bound must hold not only at $X=G$. It must
hold for all $X$ with $|X-G|\le e^{-\rho N}$, and the exponent $\rho$ enters the final bound linearly.
\item \emph{Uniform nonvanishing.} One needs $Q_p(a/b)\ne0$ for primes $p$ with $p^2\nmid b$, in place of
$Q_p(G)\ne0$.
\item \emph{Arithmetic.} The integer $b^{\deg Q_N}Q_N(a/b)$ is nonzero, so it is at least $1$ in absolute
value. The smaller the denominator of $\Delta_N$ and its degree, the stronger the conclusion. The prime
$p$ must be chosen not to divide $b$, and a naive choice costs a factor.
\end{enumerate}
Let $d$ be the degree of $Q_N$ divided by $N$ and $\eta$ the decay rate of $Q_N$ normalised by $N^2$.
Then the conclusion is $\mu(G)\le d\rho/\eta$ (Proposition~\ref{prop:criterion}). The work is to make $d$
and $\rho$ small and $\eta$ large. Since $\eta$ is a small difference of two large numbers, the gains in
$\eta$ matter most.

\subsection{New ideas}\label{sec:intro-ideas}
\begin{enumerate}[label=\arabic*.,wide]
\item \emph{All parameters at once.} The determinant has five integer parameters
$(n_0,b_0,q_0,g_0,h_0)$, which we call a \emph{design}. Every argument below is carried out for an
arbitrary design, with the parameters as symbols. This covers the degree, the denominators at all
primes, the rank deficiencies, the nonvanishing and the estimates at the real place. A design is then
data. We use $(58,9,3,5,3)$, so that $n=58N$ and the degree is $53N$.

\item \emph{Rank of the reduction modulo $p$, in closed form.} Let $\nu_N(p)$ be the dimension of the left
kernel of the reduction modulo $p$ of $p^2$ times the matrix of rational parts. At a prime $p$ the
valuation of $\Delta_N$ gains $\nu_N(p)$. We prove four lower bounds for $\nu_N(p)$
(Theorem~\ref{thm:content}). Each is a piecewise linear function of $p$, with coefficients linear in the
five parameters, on its own window of primes. They come from antisymmetric combinations of rows, a
reflection of Chebyshev polynomials modulo $p$, and the Frobenius relation $(1-X)^p=1-X^p$. All four are
proved from one digit formula for the entries modulo $p$.

\item \emph{A counting inequality for the nullity.} The four proofs share one shape. A space $\cV$ of row
combinations is given, with a few linear conditions on $\cV$, and the conditions force all column
responses to vanish. We prove this once (Theorem~\ref{thm:count}). Then $\nu_N(p)\ge\dim\cV-\rank\Phi$,
where $\Phi$ is the map of conditions. A congruence from Lucas's theorem lets the inequality apply when
the column filter is longer than the prime. Further families of kernel vectors are instances of it.

\item \emph{A fixed rational subspace with stored kernel vectors.} In one form of the inequality the
space is a tensor product of polynomial coordinates with a fixed subspace $V\subset\Q^s$
(Theorem~\ref{thm:profile}). The required ranks are certified by integer vectors in the kernels of fixed
matrices, stored once. This treats $34$ windows of primes between $3N$ and $16.2N$.

\item \emph{Paying for the perturbation column by column.} The smallness near $G$ comes from a Gauss-type
quadrature formula for the coefficient of $G$, with $53N$ nodes. A perturbed determinant is a sum of
determinants in which some columns are replaced by point values. We bound each replaced column separately.
A binomial bound for generalized Vandermonde determinants (Lemma~\ref{lem:binomvand}) and a contour bound
for the weight of each column (Lemma~\ref{lem:keyhole}) give costs that peak at different columns. Their
sum peaks lower than the sum of their peaks.

\item \emph{A signed bound for the determinant of the row functions.} OpenAI's integrand contains a
determinant $\det\mathsf L$ that the energy method bounds only by $6^n$. We bound it on the configurations
that matter (Section~\ref{sec:signed}). A Pick matrix argument and a Schur complement reduce
$|\det\mathsf L|$ to the block of the negative nodes, at a cost of $11$ per positive node. The key input is
a contour bound for a quadratic form on the model space of a Blaschke product (Lemma~\ref{lem:T}).
Heine's formula and the Christoffel formula write the negative block as a ratio of two log-gas
integrals. We bound the numerator from above by completing squares in Chebyshev moments, and the
denominator from below by Jensen's inequality. An inequality for the logarithmic energy of signed
measures (Theorem~\ref{thm:reflection}) keeps the combined bound concave. A separate argument treats
positive nodes close to $1$.

\item \emph{Near-optimal configurations.} The signed bound needs hypotheses on the configuration. They hold
on every configuration whose energy is within the gain of the maximum. On the other configurations OpenAI's
bound is already strong enough. So the improved bound holds everywhere.

\item \emph{Trial potentials with complex poles, interpolated in the number of fixed columns.} The energy
estimates use trial sequences whose generating functions are logarithmic potentials of point charges near
the edges of the support of the extremal measure. The signed bound uses three trials, at fractions
$0$, $1/200$ and $1/16$ of fixed columns. Between them the fields are affine in the fraction, and the
remaining constant is an explicit quadratic.

\item \emph{A criterion with a free lower end of the window.} In the choice of the prime, the window of
candidates is $(V,U]$ with $U=\kappa\log b$ and $V=(\kappa-\omega)\log b$. We prove the criterion for every
$\omega\in(0,1]$ with the hypothesis $\kappa>2\omega$ (Proposition~\ref{prop:criterion}). We use
$\omega=3/5$.
\end{enumerate}

\subsection{The constants}
For the design $(58,9,3,5,3)$ the degree of $Q_N$ is $d=53N$. The perturbation exponent is
$\rho=\Rho$. The decay rate at the real place is $\TDelta\ldots$, and the denominators cost
$\mathsf A=\Aden\ldots$, after a gain of $\int\kappa\ge\ChiLean$ from the rank deficiencies. The signed
bound contributes $\Gain$ to the decay rate. The criterion uses $\eta=\EtaTot$, $\omega=3/5$,
$\delta=2.9\cdot10^{-7}$ and $\kappa=410/\rho=\Kappa\ldots$ Its last condition reads $B>\Bx\ldots$, and
$410$ is the least integer above it. Decimals are truncated. The constants of the formal proof are the
exact rational numbers.

\subsection{Organization and notation}
Section~\ref{sec:overview} sets up designs and states the rank deficiencies, the valuations, the
denominators, the bound at the real place and the criterion. Section~\ref{sec:setup} recalls OpenAI's
construction for a design and proves the degree, the denominators and the nonvanishing.
Section~\ref{sec:content} proves the rank deficiencies. Section~\ref{sec:real} reduces the bound near
$G$ to determinants with fixed columns and recalls the energy method. Section~\ref{sec:signed} proves the
signed bound. Section~\ref{sec:proof} proves Theorem~\ref{thm:main}. Section~\ref{sec:remarks} makes
remarks on the method, and Section~\ref{sec:formal} describes the formalization.

We use the notation of OpenAI's paper \cite{OpenAI2026} for the construction and refer to it for details.
Citations of the form \cite[Prop.~2.1]{OpenAI2026} point to its numbered statements. We write $\log$ for
the natural logarithm, $v_p$ for the $p$-adic valuation and $\theta(x)=\sum_{p\le x}\log p$ for
Chebyshev's function. $T_k$ and $U_k$ are the Chebyshev polynomials of the first and second kind, with
$U_{-1}=0$, $T_{-k}=T_k$ and $U_{-k-1}=-U_{k-1}$. We write $[X^k]F$ or $F[k]$ for the coefficient of $X^k$
in a polynomial $F$, $y_+=\max(y,0)$, and $\operatorname{rev}_EF(X)=X^EF(1/X)$ for the reversal of a
polynomial of degree at most $E$.

\section{Designs, rank deficiencies and the exponent}\label{sec:overview}

\subsection{Designs}
A \emph{design} is a quintuple $(n_0,b_0,q_0,g_0,h_0)$ of integers with $n_0,b_0,q_0,h_0\ge1$,
$g_0\ge0$ and $2g_0<n_0$. For $N\ge1$ we write
\begin{gather*}
 n=n_0N,\quad b=b_0N,\quad q=q_0N,\quad g=g_0N,\quad h=h_0N,\\
 L=n+b+q,\quad C=L+g,\quad H=C+h,\quad A=b+q+2g ,
\end{gather*}
and $d_0=n_0-g_0$. OpenAI's choice is the design $(48,7,4,4,2)$. We work throughout with
\begin{gather*}
 D=(58,9,3,5,3):\quad n=58N,\quad b=9N,\quad q=3N,\quad g=5N,\quad h=3N,\\
 L=70N,\quad C=75N,\quad H=78N,\quad A=22N,
\end{gather*}
and $d_0=53$. Section~\ref{sec:setup} recalls the construction of OpenAI's determinant for a design. It
yields an $n\times n$ matrix $\cA_N=\cA^0_N+G\cL_N$ with rational matrices $\cA^0_N$ and $\cL_N$. The
polynomial $\Delta_N(X)=\det(\cA^0_N+X\cL_N)$ has degree at most $n-g$ (Lemma~\ref{lem:degree}). The search
that led to $D$ is not part of the proof.

\subsection{Rank deficiencies}
For an odd prime $p$ with $H<p^2$ the entries of $p^2\cA^0_N$ are $p$-integral (Lemma~\ref{lem:digits}).
Let $\cA^0_{N,p}$ be its reduction modulo $p$, and let $\nu_N(p)$ be the dimension of its left kernel, that
is, $n$ minus its rank. Write $\Lambda=g+q+h$ and $e=(p+1)/2-q-h$.

\begin{theorem}\label{thm:content}
Let $D$ be a design, $N\ge1$, and $p$ an odd prime with $H<p^2$. Then $\nu_N(p)$ is at least each of the
following quantities whenever the stated conditions hold. Negative values are replaced by $0$.
\begin{enumerate}[label=(\alph*)]
\item (Theorem~\ref{thm:T}) If $2g<p$, $H\le4p$ and $D_T=\min(g,\,n-1-g-2p,\,C-3p)\ge1$:
\[
 D_T-(n+q+g+h-3p)_+ .
\]
\item (Theorem~\ref{thm:U}) If $2h\le p+1$, $L\le3p\le H$, $n+q+g+h\le3p$ and $2p\le n+q$:
\[
 \min(n,\,2p+2g+1)-g-2p-1 .
\]
\item (Theorem~\ref{thm:U6}) If $e\ge0$, $n\le4p+2g+1$, $L\le6p$ and $b+q\le p+\mu_0$, with
$\mu_0=\max(q,H+q-6p)$:
\[
 \min\Bigl(\frac{p-1}2-g,\ n+b+q+h-\frac{11p+1}2\Bigr).
\]
\item (Theorem~\ref{thm:S}) If $e\ge0$ and $n\le4p+2g+1$: $m_1-m_0-e$, with
\[
 m_0=\max(b+q,\,L+\Lambda-4p),\qquad m_1=\min(p,\,L-3p,\,b+q+2g+1).
\]
\end{enumerate}
\end{theorem}

These are the Lean theorems \lean{t\_rank}, \lean{u\_rank}, \lean{u6\_rank\_closed} and \lean{s\_rank},
proved for every admissible design. For $D$ they apply for $12N<p\le26N$, where they give most of the gain
(Section~\ref{ssec:kappa}).

Below $12N$ the closed formulas give nothing. Further kernel vectors come from a single counting
inequality, of which the proofs of Theorem~\ref{thm:content} are special cases. For $c\in\F_p^n$ and a raw
column $j\in[b,L)$ let $R_c(j)$ be the reduction modulo $p$ of $p^2\sum_rc_rR(r,j)$, with $R$ as in
\eqref{eq:rational-part}. These are the responses of Lemma~\ref{lem:digits}.

\begin{theorem}\label{thm:count}
Let $D$ be a design, $N\ge1$, $p$ an odd prime with $H<p^2$, and write $q=ap+r$ with $0\le r<p$. Let
$\cV$ be a finite-dimensional vector space over $\F_p$, $\iota\colon\cV\to\F_p^n$ an injective linear map and
$\Phi$ a linear map on $\cV$. Suppose that for every $x\in\ker\Phi$ there is a polynomial $G_x\in\F_p[X]$ of
degree less than $r$ with
\[
 \sum_{v=0}^a(-1)^{a-v}\binom av\,R_{\iota(x)}(j+vp)=G_x(j)\qquad(b\le j<L-ap).
\]
Then $\nu_N(p)\ge\dim\cV-\operatorname{rank}\Phi$.
\end{theorem}

For $q<p$ the hypothesis says that the conditions $\Phi$ force all raw responses to vanish. The case used
for the counting is the \emph{profile form}. There the space is a tensor product of a space of polynomial
coordinates with a fixed rational subspace, and the conditions are jets composed with fixed linear maps.

\begin{theorem}\label{thm:profile}
Let $v_1,\dots,v_d\in\Z^s$ span a subspace $V\subset\Q^s$ of dimension $d$. For $j$ in a finite set $J$ let
$B_j$ be an integer matrix with $s$ columns and $r_j\le d$ an integer such that there are $d-r_j$ linearly
independent vectors $k\in\Z^d$ with $B_j\bigl(\sum_ak_av_a\bigr)=0$. Let $\mathcal E\ne0$ be the product of a
nonzero $d\times d$ minor of $(v_1,\dots,v_d)$ and, for each $j$, of a nonzero maximal minor of these kernel
vectors. Let $p\nmid\mathcal E$ be an odd prime with $H<p^2$, write $q=ap+r$ as above, let $V_p=\F_p^d$ be the
space of coordinates of $V$ modulo $p$, and let $m$ and $w_j$ ($j\in J$) be integers $\ge0$. Suppose that there
are linear maps $\varphi_{j,k}\colon V_p^m\to V_p$ ($k<w_j$) and $\iota\colon V_p^m\to\F_p^n$ such that, on the joint
kernel of the maps $B_j\circ\varphi_{j,k}$, the map $\iota$ is injective and satisfies the hypothesis of
Theorem~\ref{thm:count}. Then
\[
 \nu_N(p)\ge m\,d-\sum_{j\in J}w_j\,r_j .
\]
\end{theorem}

In the applications $m$ and the $w_j$ are affine functions of $N$ and $p$ on a window of primes, while $V$
and the $B_j$ are fixed. So the bound is linear on the window. The data to be stored are the vectors $v_a$,
the matrices $B_j$ and the kernel vectors. For the design $D$ these bounds, the closed formulas and the
further families of Section~\ref{sec:content} combine into the following.

\begin{theorem}\label{thm:profile-D}
For the design $D$ there are $102$ intervals $(a_i,b_i]$, with $16/7=a_1<b_1\le a_2<\dots<b_{102}=53/2$, and a
function $\kappa$, linear on each interval and $0$ outside their union, with the following property. For all
large $N$ and every prime $p$ with $a_iN<p\le b_iN$ for some $i$,
\[
 \nu_N(p)\ge N\kappa(p/N)-2 .
\]
On each interval the bound is given by one of the families of Section~\ref{sec:content}. The area of
$\kappa$ is a rational number,
\[
 \int_0^\infty\kappa(x)\,dx=\frac{36235515874601938601596598}{1642609663271115992578125}=\Chi\ldots\ \ge\ \ChiLean .
\]
\end{theorem}

The contributions by ranges of $x=p/N$ are as follows. The third column counts the intervals on which the
bound is that of Theorem~\ref{thm:profile} with stored data.
\begin{center}
\begin{tabular}{@{}lrrl@{}}
\toprule
range of $x$ & intervals & with stored data & area \\
\midrule
$(16/7,\,3]$ & $11$ & $0$ & $0.06259$ \\
$(3,\,6]$ & $46$ & $23$ & $0.25136$ \\
$(6,\,12]$ & $27$ & $9$ & $0.71965$ \\
$(12,\,53/2]$ & $18$ & $2$ & $21.02610$ \\
\midrule
total & $102$ & $34$ & $\Chi$ \\
\bottomrule
\end{tabular}
\end{center}
The intervals in $(12,53/2]$ are listed in Section~\ref{ssec:kappa}. The full list is part of the
formalization. On each interval one bound is used. The families come from different arguments. Except for one
pair (Section~\ref{ssec:families}) their kernels are not shown to be independent, so at a prime the bounds are
not added.

\subsection{Valuations, denominators and the exponent}
The rank deficiencies enter through the following statement.

\begin{proposition}\label{prop:valuation}
Let $p$ be an odd prime with $H<p^2$ and let $z\in\Q$ have denominator prime to $p$. Then
$p^2(\cA^0_N+z\cL_N)$ has $p$-integral entries, its reduction modulo $p$ equals $\cA^0_{N,p}$, and
\[
 v_p\bigl(\Delta_N(z)\bigr)\ge\nu_N(p)-2n .
\]
\end{proposition}

At the other primes there are the valuation bounds of OpenAI's method (Section~\ref{sec:setup}). At $2$
there is a bound from the expansion into minors. Together they give the denominators.

\begin{proposition}\label{prop:denominators}
For the design $D$ there are positive integers $D_N$ with the following properties. $D_N\Delta_N(i)\in\Z$ for
every integer $i$, and $q^{2n}\mid D_N$ for every prime $N<q\le2N$. For every $\varepsilon>0$ and all large $N$
\begin{gather*}
 \log D_N\le(\mathsf A+\varepsilon)N^2,\\
 \mathsf A=\frac{2673}8\log2+116+\frac{24269}4-\int_0^\infty\kappa(x)\,dx=\Aden\ldots
\end{gather*}
\end{proposition}

Here $\frac{2673}8\log2$ is the contribution of the prime $2$. The number $116+\frac{24269}4$ is the area of
the piecewise linear function that bounds the poles at the odd primes (Section~\ref{sec:setup}). The last
term is the gain of Theorem~\ref{thm:profile-D}.

At the real place we prove the following in Sections~\ref{sec:real} and~\ref{sec:signed}. Put
\begin{equation}\label{eq:rho-def}
 \rho=n_0c_2+\Psi_0=\Rho,\qquad c_2=\Ctwo,\quad \Psi_0=\PsiZero .
\end{equation}
Here $c_2$ is the energy cost of one fixed column and $\Psi_0$ bounds the logarithmic size of the weight
and the binomial factor of one column, per unit of $N$.

\begin{theorem}\label{thm:real}
Let $U_0=U+\Gain/n_0^2=\Uzero\ldots$, with $U=\Utarget$ and $n_0=58$. For every $\varepsilon>0$, all large
$N$ and all real $x$ with $|x-G|\le e^{-\rho N}$,
\[
 |\Delta_N(x)|\le\exp\bigl(-(n_0^2(U_0-\tfrac12\log2)-\varepsilon)N^2\bigr).
\]
\end{theorem}

The number $n_0^2(U-\frac12\log2)=\TDeltaOld\ldots$ is the decay rate that OpenAI's energy method gives
for the design $D$. The signed bound adds $\Gain$, so $n_0^2(U_0-\frac12\log2)=\TDelta\ldots$ Since
$\mathsf A=\Aden\ldots$, the polynomial $Q_N=(d_0N)!\,D_N\Delta_N$ has integer coefficients, degree at
most $d_0N$, and
\[
 |Q_N(x)|\le e^{-\eta N^2}\qquad(|x-G|\le e^{-\rho N}),\qquad \eta=\EtaTot .
\]
The exact difference is $\EtaExact\ldots$, and $\eta$ is a rational number just below it. The remaining
properties of Section~\ref{sec:intro-need} are the nonvanishing of $Q_p(a/b)$ for $p^2\nmid b$
(Proposition~\ref{prop:nonvanishing}) and the divisibility of the $i$-th coefficient of $Q_N$ by $q^{2i}$
for $N<q\le2N$ (Lemma~\ref{lem:recycling}). They are used to choose the prime. The criterion is the
following.

\begin{proposition}\label{prop:criterion}
Let $\xi\in\R$, $d\in\N$, and let $\eta,\rho,\kappa,\delta,\omega>0$ satisfy
\[
 \omega\le1,\quad 2\omega<\kappa,\quad \eta<2d,\quad 1<2\bigl(\omega-\delta(2\kappa-\omega)\bigr),\quad
 d+4d\delta\kappa<\eta\kappa .
\]
Let $Q_p\in\Z[X]$ ($p$ prime) satisfy the following for all large $p$. $\deg Q_p\le dp$. $Q_p(a/b)\ne0$
whenever $a\in\Z$, $b\ge1$ and $p^2\nmid b$. For all primes $p<q\le2p$ and $i\le dp$, $q^{2i}$ divides the
$i$-th coefficient of $Q_p$. Finally $|Q_p(x)|\le e^{-\eta p^2}$ for $|x-\xi|\le e^{-\rho p}$. Then
\[
 |\xi-a/b|>b^{-\rho\kappa}
\]
for all integers $a$ and all sufficiently large $b$.
\end{proposition}

The conditions say, in the order of the proof in Section~\ref{sec:proof}, three things. The window of primes
contains a prime not dividing $b$. The primes of the window above the chosen one are recycled. The final
inequality is contradictory. For the design $D$ one takes $d=53$, $\omega=3/5$, $\delta=2.9\cdot10^{-7}$
and $\kappa=B/\rho$ with $B=410$. The last condition then reads $B>x:=d\rho/(\eta-4d\delta)=\Bx\ldots$

\subsection{A reading of the exponent}
For a fixed design the bound is $d\rho/\eta$ in the limit. Its three ingredients come from three
different places. The degree $d=n-g$ comes from the construction. The perturbation exponent $\rho$ comes
from the quadrature and the energy cost of a column. The decay rate $\eta$ is the difference of the real
place and the denominators, $\TDelta$ and $\Aden$. Both large numbers are improved here: the rank
deficiencies lower the second by $22.05$, and the signed bound raises the first by $\Gain$. Without either,
$\eta$ would be $\EtaInf$.

\section{OpenAI's construction for a design}\label{sec:setup}

We recall the construction of \cite{OpenAI2026} in the generality in which we use it, and we prove for an
arbitrary design the degree bound, the bounds for the denominators at the finite places and the
nonvanishing at the primes $N=p$. Carreira \cite{Carreira2026} checks that the argument of OpenAI goes
through for designs satisfying a few conditions. We need the same for the finite places and, in addition,
for the entries modulo $p^2$.

\subsection{The determinant}
Let $D=(n_0,b_0,q_0,g_0,h_0)$ be a design, with the notation of Section~\ref{sec:overview}. For $0\le r<n$
the row polynomials are
\begin{equation}\label{eq:rows}
 P_r(t)=(1-t)^ht^{C-1}T_{|r-g|}(1/t),\qquad \deg P_r<H,
\end{equation}
together with companion polynomials $D_r$ built from $U_{|r-g|-1}$ \cite[(3)]{OpenAI2026}. Thus
$P_r$ depends only on $|r-g|$. Put $f(t)=\sqrt{1-t^2}$ and $d\mu(t)=|t|\,dt/f(t)$ on $(-1,1)$, and
\[
 A_r(t)=P_r(t)-\frac32\,\mathbf 1_{\{t>0\}}\frac{f(t)}{t}D_r(t),\qquad
 \psi_k(s)=s^{b+k}(1-s)^q .
\]
The determinant studied by OpenAI is that of the $n\times n$ matrix
\begin{equation}\label{eq:matrix}
 \cA_N(r,k)=\int_{-1}^1\int_0^1 A_r(t)\,\frac{\psi_k(s)}{1-ts}\,ds\,d\mu(t)
 \qquad(0\le r,k<n).
\end{equation}
The factor $(1-s)^q$ in $\psi_k$ amounts to a finite difference of order $q$ applied to the columns.
The entries are rational linear combinations of $1$ and $G$: there are rational matrices $\cA^0_N$
and $\cL_N$ with
\begin{equation}\label{eq:L}
 \cA_N=\cA^0_N+G\,\cL_N,\qquad
 \cL_N(r,k)=\sum_{v=0}^{q}(-1)^v\binom qv\sum_{i}[t^i]P_r\cdot4\,c_{(i-b-k-v)/2},
\end{equation}
where $c_l=4^{-l}\binom{2l}{l}$ for integers $l\ge0$ and $c_z=0$ otherwise
\cite[Prop.~2.1 and (20)]{OpenAI2026}. The rational part is
\begin{equation}\label{eq:rational-part}
 \begin{gathered}
 \cA^0_N(r,k)=\sum_{v=0}^q(-1)^v\binom qv R(r,b+k+v),\\
 R(r,j)=\sum_{i<H}P_r[i]\,M(i,j)-\frac32\sum_{i<H}D_r[i]\,Z(i,j)
 \end{gathered}
\end{equation}
with rational numbers $M(i,j)$ (moments of $\mu$ against the Cauchy kernel) and $Z(i,j)$, where
$Z(i,i)=-\sum_{k\le i}k^{-2}$ and $Z(i,j)=(H_i-H_j)/(i-j)$ for $i\ne j$, $H_i=\sum_{k\le i}1/k$
\cite[Prop.~2.1]{OpenAI2026}. The design enters \eqref{eq:rational-part} only through the integer
coefficients $P_r[i],D_r[i]$ and the range of $j$, which is $[b,L)$. Put
$\Delta_N(X)=\det(\cA^0_N+X\cL_N)\in\Q[X]$, so that $\Delta_N(G)=\det\cA_N$.

\subsection{The degree}
\begin{lemma}\label{lem:degree}
$\deg\Delta_N\le n-g$.
\end{lemma}

\begin{proof}
By \eqref{eq:L} the row $r$ of $\cL_N$ depends only on $P_r$, and hence only on $|r-g|$. So the rows
$r$ and $2g-r$ of $\cL_N$ agree for $0\le r<g$, and $2g-r\le n-1$ because $2g<n$. Subtracting row $r$ from
row $2g-r$ for each $r<g$ does not change the determinant and makes $g$ rows of the new coefficient of
$X$ vanish. A determinant $\det(B+XL')$ in which $L'$ has $n-g$ nonzero rows has degree at most $n-g$ in $X$.
\end{proof}

\subsection{Rank deficiency and valuation}
We use repeatedly the following elementary fact.

\begin{lemma}\label{lem:rank-valuation}
Let $M$ be an $n\times n$ matrix with entries in $\Z_{(p)}$ whose reduction modulo $p$ has rank at most
$n-e$. Then $v_p(\det M)\ge e$.
\end{lemma}

\begin{proof}
Factor the reduction through $\F_p^{n-e}$ and lift the two factors to integral matrices $U,V$, so that
$M=UV+pE$ with $E$ integral at $p$. Expanding the determinant multilinearly in the columns, every term
that uses more than $n-e$ columns of $UV$ vanishes, and every other term has at least $e$ factors $p$.
\end{proof}

\subsection{The digit formula}
The reduction of $p^2\cA^0_N$ modulo $p$ is described by the digits in base $p$ of the coefficients of the
row polynomials. This is the one place where the arithmetic of the moments enters. It is design-free.
For an odd prime $p$ and $c\in\F_p^n$ put
\[
 P_c=\sum_rc_rP_r,\qquad D_c=\sum_rc_rD_r,\qquad W=(1-X^2)^{(p-1)/2},
\]
(reading the $P_r,D_r$ as polynomials in $X$), and let $F[k]$ denote the coefficient of $X^k$.

\begin{lemma}[digit formula]\label{lem:digits}
Let $p$ be an odd prime with $H<p^2$, let $0\le j<H$, and write $j=Jp+\ell$ with $0\le\ell<p$. Then
$p^2R(r,j)$ is $p$-integral for every row $r$, and for every $c\in\Z^n$
\[
 p^2\sum_rc_rR(r,j)\equiv\sum_{u\ge0}\Bigl(\alpha_{u,J}\,(XWP_c)[up+\ell]+\beta_{u,J}\,D_c[up+\ell]\Bigr)
 \pmod p ,
\]
where $\alpha_{u,J}=M(u-1,J)$ and $\beta_{u,J}=-\tfrac32Z(u,J)$ are $p$-integral numbers that do not depend on
$c$, on $\ell$ or on the design.
\end{lemma}

\begin{proof}
Two facts about the numbers $M$ and $Z$ are needed, for $0\le i,j<p^2$. Write $j=Jp+\ell$ and let
$\varrho\in[0,p)$ be the residue of $j-i-1$ modulo $p$, $u=(i+1+\varrho-\ell)/p$. First, $p^2M(i,j)$ is
$p$-integral and $p^2M(i,j)\equiv W[\varrho]\,M(u-1,J)\pmod p$. This is the digit reduction of the moments in
\cite[\S3]{OpenAI2026}, obtained from the $p$-adic analysis of their recursions, and it does not involve a design.
Second, $p^2Z(i,j)$ is $p$-integral and is congruent to $Z(\lfloor i/p\rfloor,J)$ if $i\equiv j\pmod p$ and to
$0$ otherwise. For the second, write $i'=\lfloor i/p\rfloor$. Since $i<p^2$, $H_i=p^{-1}H_{i'}+S_i$ and
$H^{(2)}_i=p^{-2}H^{(2)}_{i'}+S'_i$ with $S_i,S'_i$ $p$-integral, which gives the case $i=j$. If $i\not\equiv j$,
then $p^2Z(i,j)=p\cdot pZ(i,j)\equiv0$. If $i\equiv j$, $i\ne j$, then $i-j=(i'-J)p$, and the integers
$k\in(j,i]$ prime to $p$ form $i'-J$ complete systems of residues modulo $p$, so that
$\sum_k1/k\equiv(i'-J)\sum_{r=1}^{p-1}1/r\equiv0\pmod p$. Hence $H_i-H_j\equiv(H_{i'}-H_J)/p$ and
$p^2Z(i,j)\equiv Z(i',J)$.

Now use \eqref{eq:rational-part}. By the first fact the part $\sum_iF[i]\,p^2M(i,j)$ with $F=P_c$ is
$\sum_iF[i]\,W[\varrho(i)]\,M(u(i)-1,J)$. For fixed $u$ the indices with $u(i)=u$ are those for which
$\varrho=up+\ell-1-i$ lies in $[0,p-1]$, and $W$ has degree $p-1$, so the sum over these $i$ is
$(XFW)[up+\ell]$. This gives the $\alpha$-terms. By the second fact the part with $D_c$ is
$\sum_{i\equiv\ell}D_c[i]\,Z(\lfloor i/p\rfloor,J)=\sum_uZ(u,J)D_c[up+\ell]$, which gives the $\beta$-terms.
The $p$-integrality of $M(u-1,J)$ and $Z(u,J)$ for $u,J<p$ is part of the first and second facts.
\end{proof}

For the filtered columns we use the following form. Let $0\le k<n$, $s=b+k$, $J=\lfloor s/p\rfloor$ and
$\mu=(s\bmod p)+q$. If $\mu<p$, the $q+1$ raw columns $s,\dots,s+q$ that enter column $k$ have the same
digit $J$ and residues $(s\bmod p)+v<p$. Since
$\sum_{v=0}^q(-1)^v\binom qvF[t+v]=[X^{t+q}]\bigl((X-1)^qF\bigr)$ for every polynomial $F$,
Lemma~\ref{lem:digits} gives
\begin{equation}\label{eq:filtered-digits}
 p^2\sum_rc_r\cA^0_N(r,k)\equiv\sum_{u\ge0}\Bigl(\alpha_{u,J}\,Q_c[up+\mu]+\beta_{u,J}\,\widehat D_c[up+\mu]\Bigr),
\end{equation}
where $Q_c=(X-1)^q\,XWP_c$ and $\widehat D_c=(X-1)^qD_c$. A combination $c$ whose responses vanish on all
raw columns in $[b,L)$ lies in the left kernel of $\cA^0_{N,p}$, and a combination whose responses on the
columns outside a set vanish and which kills the set by \eqref{eq:filtered-digits} does as well.

\subsection{Integrality and the valuation at $p$}
\begin{proof}[Proof of Proposition~\ref{prop:valuation}]
By Lemma~\ref{lem:digits} the entries of $p^2\cA^0_N$ are $p$-integral, since the filter is an integer
combination of raw columns with $j<L\le H$. The entries of $\cL_N$ are $p$-integral for odd $p$, since
$P_r[i]\in\Z$ and $c_l\in\Z[1/2]$, so $p^2\cL_N\equiv0\pmod{p^2}$ and, for $z$ with denominator prime to
$p$, $p^2(\cA^0_N+z\cL_N)$ has $p$-integral entries with the same reduction as $p^2\cA^0_N$. By
Lemma~\ref{lem:rank-valuation} the determinant of the latter has valuation at least $\nu_N(p)$, and
$\det\bigl(p^2(\cA^0_N+z\cL_N)\bigr)=p^{2n}\Delta_N(z)$.
\end{proof}

\subsection{The denominators}
Let $p$ be an odd prime with $H<p^2$ and $p\le H$. The method of \cite[\S3]{OpenAI2026}, applied to a
design, bounds from below the valuation of $\Delta_N(z)$ for every $z$ with denominator prime to $p$ by
$-\ell_N(p)$, where $\ell_N(p)$ counts the columns whose entries have a pole at $p$, once the coefficients of
the row polynomials are grouped by their digits in base $p$ and the column filter is taken into account. It
is given by two formulas, in which all differences are truncated at $0$. For $2p\le H$,
\[
 \ell_N(p)=2n-\Bigl(\min\bigl(p,\,L-p,\,A\bigr)-\max\bigl(b,\,H-2p\bigr)-q\Bigr),
\]
and for $H<2p\le2H$, with $R=(L-b)-(\min(A,L-p)-b)-(\min(p,L)-\max(b,H-p))$,
$P'=\min(A,L-p)-b$ and $E'=\min(b,H-p)-(L-p)$,
\[
 \ell_N(p)=\min\bigl(2n,\ n+R,\ 2R+P'+E'\bigr).
\]
These are the exponents of the generic valuation theorems in the formalization. For the primes that
interest us they are the same as in OpenAI's paper for its design. For the design $D$ the function
$\ell_N(p)=Nf(p/N)$ is piecewise linear in $x=p/N$ ($\lfloor\cdot\rfloor$ rounding aside):
\[
 f(x)=\begin{cases}
 116&x\le59/2,\\
 175-2x&59/2\le x\le69/2,\\
 106&69/2\le x\le39,\\
 184-2x&39\le x\le81/2,\\
 265-4x&81/2\le x\le48,\\
 217-3x&48\le x\le69,\\
 148-2x&69\le x\le70,\\
 78-x&70\le x\le78,
 \end{cases}
\]
and $f=0$ for $x\ge78$, since for $p>H$ the matrix $\cA^0_N+z\cL_N$ is $p$-integral. Its integral over
$[1,78]$ is $24269/4$, and $f=116=2\cdot58$ on $[0,1]$. Note that $f(x)=2n_0$ for $x\le59/2$, which
contains all the primes in $(N,2N]$ and, in particular, the whole range $(16/7,53/2]$ of the rank
deficiencies of Theorem~\ref{thm:profile-D}. There the bound of Proposition~\ref{prop:valuation} improves
$-\ell_N(p)=-2n$ to $\nu_N(p)-2n$.

The primes $p\le\sqrt H$ are not covered by these bounds. For them we use the cruder bound
$6n\log_pH$ for the exponent, which contributes $O(\sqrt N\,N\log N)=o(N^2)$ in total.

\subsection{The prime $2$}
At $p=2$ the entries are not integral and the bound comes from the expansion of the determinant into
minors, in the manner of \cite[\S3]{OpenAI2026}. Its exponent is governed by a quadratic form in the
numbers $m,\ell$ of columns of two exceptional types, which is
\[
 \tfrac32m^2+\tfrac12\ell^2-(H+C-b)\,m-C\,\ell+Cn ,
\]
and the following identity shows that it is at least $-\mathsf A_2$ whenever $m,\ell\ge0$,
$m+\ell\le n$ and $C\ge n$, where
\begin{equation}\label{eq:A2}
 \mathsf A_2=\frac{(n+H-b)^2}8-\frac{n^2}2 .
\end{equation}
With $s=n-m-l$ and $m_*=(n+H-b)/4$,
\[
 \tfrac32m^2+\tfrac12\ell^2-(H+C-b)m-C\ell+Cn+\mathsf A_2
 =2(m-m_*)^2+s\,(C-n+m)+\tfrac12s^2\ge0 .
\]
This gives a $2$-adic bound $v_2(\Delta_N(i))\ge-\mathsf A_2-O(N\log N)$, uniform in the integer $i$. For $D$,
$\mathsf A_2=(58+78-9)^2/8-58^2/2=2673/8$ in units of $N^2$.

\subsection{The common denominator}
\begin{proof}[Proof of Proposition~\ref{prop:denominators}]
Let $\kappa_N(p)$ be the largest of the lower bounds for $\nu_N(p)$ of Section~\ref{sec:content} that apply at
$p$ (and $0$ if none does), for primes $p\le H$ with $H<p^2$, and put $e_N(p)=\min\bigl(\ell_N(p),\,2n-\kappa_N(p)\bigr)$. For the
primes with $p^2\le H$ use the crude bound above, for $p=2$ the $2$-adic bound, and let $D_N$ be the
product of $p^{e_N(p)}$ over $p\le H$. By Proposition~\ref{prop:valuation} and the bounds of this section,
$v_p(\Delta_N(i))\ge-e_N(p)$ for every integer $i\ge0$ and every $p$ (for $p>H$ the matrix is
$p$-integral), so $D_N\Delta_N(i)\in\Z$. A prime $q\in(N,2N]$ has $\kappa_N(q)=0$, since no window
contains it, and $\ell_N(q)=2n$, so $q^{2n}\mid D_N$ (here $N\ge79$ guarantees $q>\sqrt H$).

For the size of $D_N$, the sum of $\ell_N(p)\log p$ over the odd primes is, by the prime number theorem in
the weighted form
\begin{equation}\label{eq:weighted-pnt}
 \frac1N\sum_{p}F(p/N)\log p\longrightarrow\int_0^\infty F(x)\,dx\qquad(N\to\infty),
\end{equation}
for bounded, Riemann integrable, compactly supported $F$, equal to $N^2\int_0^{78}f(x)\,dx+o(N^2)$ with
$\int_0^{78}f=116+24269/4$. The primes in the intervals of Theorem~\ref{thm:profile-D} lie in
$(16N/7,53N/2]$ and have $\ell_N(p)=116N$, so their exponent is at most $116N-N\kappa(p/N)+2$. The sum of
$(N\kappa(p/N)-2)\log p$ over these primes is $N^2\int\kappa+o(N^2)$ by \eqref{eq:weighted-pnt}. The $2$-adic part contributes
$\mathsf A_2\log2\,N^2+o(N^2)$.
\end{proof}

The polynomial $\Delta_N$ has degree at most $n-g=d_0N$ and takes values in $D_N^{-1}\Z$ at the
integers. An integer-valued polynomial of degree $m$ has coefficients in $\frac1{m!}\Z$ (Newton's
interpolation formula), so
\begin{equation}\label{eq:QN}
 Q_N=(d_0N)!\,D_N\,\Delta_N(X)\in\Z[X],\qquad \deg Q_N\le d_0N,
\end{equation}
and $\log(d_0N)!=o(N^2)$. In this way the gain from the rank deficiencies enters through the integer
polynomial itself, and not through a separate common factor of its coefficients.

\begin{lemma}\label{lem:recycling}
For all large $N$ and every prime $N<q\le2N$, the $j$-th coefficient of $Q_N$ is divisible by $q^{2j}$.
\end{lemma}

\begin{proof}
For these primes the matrices $q^2\cA^0_N$ and $\cL_N$ are $q$-integral by Lemma~\ref{lem:digits}, and
$q^{2n}$ divides $D_N$. The identity
\[
 Q_N(X/q^2)=(d_0N)!\,\frac{D_N}{q^{2n}}\det\bigl(q^2\cA^0_N+X\cL_N\bigr)
\]
shows that $Q_N(X/q^2)$ has $q$-integral coefficients. Its $j$-th coefficient is $c_{N,j}/q^{2j}$.
\end{proof}

\subsection{Nonvanishing}
\begin{proposition}\label{prop:nonvanishing}
Let $D$ be a design and suppose that the three $n_0\times n_0$ matrices $B_\sigma$, $\sigma\in\{0,1,-1\}$,
defined below, are invertible. Then for all large primes $p$ and every rational $z$ such that $pz$ has
denominator prime to $p$, $\Delta_p(z)\ne0$. In particular $\Delta_p(a/b)\ne0$ whenever $a\in\Z$, $b\ge1$
and $p^2\nmid b$.
\end{proposition}

The matrices are those of OpenAI's reduction modulo $p$ at $N=p$, with the design in place of $48$: let
$B$ be the rational $(n_0+1)\times n_0$ matrix of the entries of the filtered rational part at $N=1$ for
the rows $0,\dots,n_0$ (the row $n_0$ is an auxiliary row), and let $B_\sigma$ have rows
$B_r+\sigma B_{r+1}$, $0\le r<n_0$. For the design $D$, $B$ has $59$ rows and $58$ columns.

\begin{proof}
Let $p>4H_0$, with $H_0=H/N=78$, so that $H=H_0p<p^2/4$, and let $p$ be prime to the numerators and denominators
of the entries of $B$ and of the three determinants. In the construction at $N=p$ write the row index
$r=pr_0+i$ with $0\le r_0<n_0$ and $0\le i<p$. Since $q=q_0p$, $C=C_0p$, $h=h_0p$, $g=g_0p$, the Frobenius
relation $(1-t)^{p}\equiv1-t^p$ expresses the row polynomials modulo $p$ through those of the design at
$N=1$ and the successor row $r_0+1$. This is where the auxiliary row enters. It is valid for rows on
both sides of $g$ because the polynomials are defined through $T_{r-g}$ and $U_{r-g-1}$ for signed
indices. The filter satisfies $s^{b_0p+k}(1-s)^{q_0p}=s^{\ell}(s^p)^{b_0+k_0}(1-s^p)^{q_0}$ in $\F_p[s]$, so
the filtered columns of one residue class use the raw columns $b_0+k_0,\dots,b_0+k_0+q_0$ of the design at
$N=1$ and never change the residue. By \cite[Lemma~3.1 and eqs.~(31), (35), (36)]{OpenAI2026}, which
need only $H<p^2/4$, the reduction of $p^2\cA^0_p$ modulo $p$ is then, after ordering the rows by $(i,r_0)$ and
the columns by $(\ell,k_0)$, the matrix $a^{\mathsf T}\otimes B_0+b^{\mathsf T}\otimes B_1$, where $a,b$ are the
fixed $p\times p$ matrices of OpenAI's identity (51), with $\det a=2^{-p(p-1)/2}$. The permutation
matrix $\Pi$ in this decomposition has the eigenvalue $0$ once and $\pm1$ each $(p-1)/2$ times, and
therefore
\[
 \det=(\det a)^{n_0}\,\det B_0\;\det(B_0+B_1)^{(p-1)/2}\det(B_0-B_1)^{(p-1)/2}\ne0
\]
in $\F_p$. Hence $p^{2n}\Delta_p(0)$ is a $p$-adic unit. For $z$ with $pz$ of denominator prime to $p$, the
difference $p^2z\cL_p=p\,(pz)\,\cL_p$ is $p$-integral with zero reduction (the entries of $\cL_p$ are
$p$-integral), so the reduction of $p^2(\cA^0_p+z\cL_p)$ is the same invertible matrix and
$v_p(\Delta_p(z))=-2n$. If $p^2\nmid b$, then $p(a/b)$ has denominator prime to $p$.
\end{proof}

For the design $D$ the invertibility of $B_0$, $B_0+B_1$ and $B_0-B_1$ is certified by exhibiting their
inverses modulo the prime $1009$, which exceeds $H_0=78$ so that every denominator of $B$ is a unit
there. The three determinants are then nonzero rational numbers. The computation modulo $1009$ certifies
the base matrices. The conclusion for all large primes is a theorem and not an extrapolation from
finitely many primes. Nothing in the proof uses the proportions of the design except $2g<n$, and $H<p^2/4$.
\section{Rank deficiencies modulo $p$}\label{sec:content}

In this section we prove Theorems~\ref{thm:content}, \ref{thm:count}, \ref{thm:profile}
and~\ref{thm:profile-D}. We begin with the four closed formulas, whose proofs show the mechanism in its
simplest form, then prove the counting inequality that underlies all of them, describe the further
families obtained from it, and finally the certified data for finitely many windows and the resulting
function $\kappa$ for the design $D$. The principle is that of the arithmetic method of Hata
and Rhin--Viola, in the form that suits a determinant: if the reduction modulo $p$ of $p^2\cA^0_N$ has
rank at most $n-e$, then by Lemma~\ref{lem:rank-valuation} and Proposition~\ref{prop:valuation} the
determinant $\Delta_N(z)$ gains a factor $p^e$ at every $z$ with denominator prime to $p$, in
particular at every integer. The gain is in the denominator, and by \eqref{eq:QN} it is passed on to
the coefficients of $Q_N$. We therefore need lower bounds for the dimension of the left kernel of
$\cA^0_{N,p}$, that is, explicit row combinations $c\in\F_p^n$ with $c^{\mathsf T}\cA^0_{N,p}=0$. The
tools are the same throughout: the identities for Chebyshev polynomials in characteristic $p$
\begin{equation}\label{eq:cheb-p}
 \begin{gathered}
 T_p=X^p,\qquad U_{p-1}=(X^2-1)^{(p-1)/2},\\
 T_{kp}=T_k(X^p),\qquad U_{kp-1}=(X^2-1)^{(p-1)/2}\,U_{k-1}(X^p)
 \end{gathered}
\end{equation}
in $\F_p[X]$ (the last two follow from $T_{kp}=T_k\circ T_p$ and $U_{kp-1}=(U_{k-1}\circ T_p)\,U_{p-1}$), the relation
$(1-X)^p=1-X^p$, and the digit formula \eqref{eq:filtered-digits}.

\subsection{Setting}
Fix a design, an integer $N\ge1$ and an odd prime $p$ with $H<p^2$. For $0\le r<n$ put $m_r=r-g$, so that
$-g\le m_r\le n-1-g$, and $n-1-g\ge g$. In the notation of \eqref{eq:rows},
\[
 P_r=(1-X)^h\operatorname{rev}_{C-1}T_{m_r},\qquad D_r=(1-X)^h\operatorname{rev}_{C-1}U_{m_r-1},
\]
where $T_{-m}=T_m$ and $U_{-m-1}=-U_{m-1}$, so that $D_g=0$. For $c\in\F_p^n$ write
$S_c=\sum_rc_r\operatorname{rev}_{C-1}T_{m_r}$ and $U_c=\sum_rc_r\operatorname{rev}_{C-1}U_{m_r-1}$, so that $P_c=(1-X)^hS_c$
and $D_c=(1-X)^hU_c$. The lowest degree of $\operatorname{rev}_{C-1}T_m$ is at least $C-1-|m|\ge a:=A=b+q+2g$, and that of
$\operatorname{rev}_{C-1}U_{m-1}$ is at least $a+1$. Put
\[
 Q_c=(X-1)^q\,XWP_c,\qquad \widehat D_c=(X-1)^qD_c,\qquad W=(1-X^2)^{(p-1)/2},
\]
so that, with $\Lambda=g+q+h$, $Q_c$ has its coefficients in $[a+1,\,L+\Lambda+p-1]$ and $\widehat D_c$ in
$[a+1,\,L+\Lambda-1]$. Recall that a vector $c$ lies in the left kernel of $\cA^0_{N,p}$ as soon as the
right side of \eqref{eq:filtered-digits} vanishes for every filtered column $k$ (for the columns with a
carry among the $q+1$ raw columns one splits the sum over the raw columns at the carry. We shall not need
this case).

\subsection{Antisymmetric combinations of rows}
\begin{theorem}\label{thm:T}
Let $2g<p$, $H\le4p$ and $D_T=\min(g,\,n-1-g-2p,\,C-3p)\ge1$. Then
$\nu_N(p)\ge D_T-(n+q+g+h-3p)_+$.
\end{theorem}

The rows used are the differences of rows placed symmetrically about the three centres $g$, $g+p$ and
$g+2p$.

\begin{proof}
For $1\le d\le D_T$ and $k=0,1,2$ let $\delta_k(d)=e_{g+kp+d}-e_{g+kp-d}\in\F_p^n$. These rows exist because
$d\le g$ and $g+2p+d\le n-1$, and the $3D_T$ vectors $\delta_k(d)$ are linearly independent, since their supports
are disjoint: the supports lie in the intervals $[kp,kp+2g]$ of row indices, which are disjoint because
$2g<p$. The row $g+kp\pm d$ has $m=kp\pm d$. By $T_{a+b}-T_{a-b}=2(X^2-1)U_{a-1}U_{b-1}$ and
$U_{a+b-1}-U_{a-b-1}=2T_aU_{b-1}$ together with \eqref{eq:cheb-p},
\begin{align*}
 T_{kp+d}-T_{kp-d}&=2(X^2-1)^{(p+1)/2}\,U_{k-1}(X^p)\,U_{d-1},\\
 U_{kp+d-1}-U_{kp-d-1}&=2\,T_k(X^p)\,U_{d-1}.
\end{align*}
Put $R_d=2(1-X)^hX^{C-d}\operatorname{rev}_{d-1}U_{d-1}$, the value of $D_{\delta_0(d)}$. Reversing at degree $C-1$,
and using that $\operatorname{rev}_{kp}(T_k(X^p))=(\operatorname{rev}_kT_k)(X^p)$, one finds, with $T_0=1$, $T_1=X$ and
$T_2=2X^2-1$,
\[
 D_{\delta_1(d)}=X^{-p}R_d,\qquad D_{\delta_2(d)}=2X^{-2p}R_d-R_d,
\]
\begin{gather*}
 XWP_{\delta_0(d)}=0,\qquad XWP_{\delta_1(d)}=X^{-p}R_d-X^pR_d,\\
 XWP_{\delta_2(d)}=2X^{-2p}R_d-2R_d .
\end{gather*}
The polynomial $R_d$ is supported in degrees $[C-d,H-1]$, which lies in $[3p,4p)$ because $d\le C-3p$ and
$H\le4p$. So at each residue $\ell$ it has a single base-$p$ digit $\varrho_\ell(d)$, in position $u=3$.
For $x=(x_0,x_1,x_2)\in\F_p^3$ the response of the combination $\sum_kx_k\delta_k(d)$ at the raw column
$Jp+\ell$ is therefore, by Lemma~\ref{lem:digits}, $\varrho_\ell(d)\,f_J(x)$ with
\[
 f_J(x)=x_0\beta_{3,J}+x_1(\alpha_{2,J}-\alpha_{4,J}+\beta_{2,J})+x_2(2\alpha_{1,J}-2\alpha_{3,J}+2\beta_{1,J}-\beta_{3,J}),
\]
which depends neither on $d$ nor on $\ell$. The raw columns with $J\ge3$ do not meet the support, since
$3p+\ell\ge C-d\ge C-g=L$. Choose $x^*\ne0$ with $f_1(x^*)=f_2(x^*)=0$. This is two linear equations in three
unknowns. For $c=\sum_ds_d\sum_kx^*_k\delta_k(d)$ the responses at the columns with $J=1,2$ vanish, and the
columns with $J=0$, which have $\ell\ge b$, respond by $f_0(x^*)\sum_ds_d\varrho_\ell(d)$, the coefficient
of $X^{3p+\ell}$ in $\sum_ds_dR_d$. These are at most $(H-3p-b)_+=(n+q+g+h-3p)_+$ linear conditions on
$s=(s_d)$. The map $s\mapsto c$ is injective, so the solutions form a space of dimension at least
$D_T-(n+q+g+h-3p)_+$ in the left kernel.
\end{proof}

\subsection{A reflection of Chebyshev polynomials}
The first ingredient of the next theorem is a reflection formula modulo $p$.

\begin{lemma}\label{lem:u3-reflect}
Let $p$ be an odd prime and $\nu=(p+1)/2$. In $\F_p[X]$, for every integer $m$,
\[
 T_{p-m}=X^pT_m-(X^2-1)^{\nu}\,U_{m-1}.
\]
\end{lemma}

\begin{proof}
Both sides satisfy the three-term recurrence in $m$, so it suffices to check $m=0$ and $m=1$. This
needs \eqref{eq:cheb-p}. For $m=0$ it is $T_p=X^p$. For $m=1$ use $T_{p+1}=XT_p+(X^2-1)U_{p-1}$ and
$T_{p-1}=2XT_p-T_{p+1}$.
\end{proof}

\begin{theorem}\label{thm:U}
Let $2h\le p+1$, $L\le3p\le H$, $n+q+g+h\le3p$ and $2p\le n+q$, and put $n^*=\min(n,2p+2g+1)$. Then
$\nu_N(p)\ge n^*-g-2p-1$.
\end{theorem}

\begin{proof}
We restrict $c$ to the rows $r<n^*$, so that $m_r\le n^*-1-g\le2p+g$. Applying Lemma~\ref{lem:u3-reflect}
with $m=m_r$ and reflecting at degree $p+C-1$ gives, for every such row,
\[
 (1-X^2)^\nu\operatorname{rev}_{C-1}U_{m-1}-X\operatorname{rev}_{C-1}T_m=-X^{p+1}\operatorname{rev}_{C-1}T_{p-m},
\]
and the right side is divisible by $X^{p+C-|p-m|}$, hence by $X^L$ since $|p-m|\le p+g=p+C-L$. Multiplying
by $(1-X)^h$ and summing, for every $c$ supported on these rows
\begin{equation}\label{eq:identityK}
 X^L\mid(1-X^2)^\nu D_c-XP_c .
\end{equation}

Let $V$ be the space of such $c$ with $(1-X^2)^\nu\mid P_c$. Since $(1-X)^h$ divides every $P_r$, this is
$\nu-h$ conditions at $X=1$ and $\nu$ conditions at $X=-1$, that is, at most $p+1-h$ linear conditions.
For $c\in V$ write $P_c=(1-X^2)^\nu Y$ and $\widetilde Q=XY$. Since $(1-X^2)^p=1-X^{2p}$,
\[
 XWP_c=(1-X^{2p})\,\widetilde Q .
\]
The polynomial $P_c$ has no coefficients below $X^{A^*}$, $A^*=C-n^*+g$, and none from $X^H$ on, so
$\widetilde Q$ is supported in $[A^*+1,\,H-p-1]$. Dividing \eqref{eq:identityK} by the unit
$(1-X^2)^\nu$ modulo $X^L$ gives $D_c[k]=\widetilde Q[k]$ for $k<L$.

Now let $H-3p\le\ell<L-2p$. The inequality $L-2p\le A^*+1$ is equivalent to $n^*\le2p+2g+1$, and
$L-2p\le p$ since $L\le3p$. By the supports, $\widetilde Q[\ell]=0$, $\widetilde Q[2p+\ell]=0$ and
$\widetilde Q[3p+\ell]=0$, and $D_c[3p+\ell]=0$ because $3p+\ell\ge H$. So with
$\varkappa=\widetilde Q[p+\ell]$ the base-$p$ digits of $XWP_c$ at the residue $\ell$ are
$(0,\varkappa,0,-\varkappa,0,\dots)$ and those of $D_c$ are $(0,\varkappa,0,0,\dots)$. Hence, by
Lemma~\ref{lem:digits}, the response at every raw column $Jp+\ell$ with such an $\ell$ is $\varkappa$ times
a constant depending on $J$ only: on $V$ these columns, which have $J\in\{0,1,2\}$, span a space of
dimension at most one for each residue $\ell$.

Put $\ell_0=H-3p$ and $\ell_1=L-2p$. The hypotheses give $\ell_0\le b\le\ell_1$. For $\ell_0\le\ell<b$ the raw
columns $j\in[b,L)$ with this residue have $J\in\{1,2\}$, and for $b\le\ell<\ell_1$ they have
$J\in\{0,1,2\}$. So the saving is $2\ell_1-\ell_0-b$, and the number of linear conditions needed to make all
responses on the $n+q=L-b$ raw columns vanish on $V$ is
$(p+1-h)+(n+q)-(2\ell_1-\ell_0-b)$. Substituting $\ell_1=L-2p$, $\ell_0=H-3p$ this is $2p+1+g$. The space of
solutions in the span of $n^*$ rows has dimension at least $n^*-g-2p-1$.
\end{proof}

\subsection{A block of columns: the digits $1$ to $5$}
\begin{theorem}\label{thm:U6}
Let $e=(p+1)/2-q-h\ge0$, $n\le4p+2g+1$, $L\le6p$, and put $\mu_0=\max(q,H+q-6p)$, $\mu_1=L-5p$. Assume
$b+q\le p+\mu_0$. Then $\nu_N(p)\ge\mu_1-\mu_0-e$, that is
\[
 \nu_N(p)\ge\min\Bigl(\frac{p-1}2-g,\ n+b+q+h-\frac{11p+1}2\Bigr).
\]
The first term is the minimum if $6p\le H$ and the second if $6p\ge H$.
\end{theorem}

Unlike Theorems~\ref{thm:T} and~\ref{thm:U}, where the conditions are imposed on the rows, the argument
concerns a block of columns. We first need an identity.

\begin{lemma}\label{lem:K6}
In $\F_p[X]$, for every integer $k$, $2X^p(X^2-1)^{(p-1)/2}T_k-(2X^{2p}-1)U_{k-1}=U_{2p-k-1}$. If
$n\le4p+2g+1$, then for every $c$, $X^L\mid 2XWS_c-(2-X^{2p})U_c$, and consequently
\begin{equation}\label{eq:K6}
 2\,Q_c[i]=2\,\widehat D_c[i]-\widehat D_c[i-2p]\qquad(i<L).
\end{equation}
\end{lemma}

\begin{proof}
For all integers $m,k$ one has $U_{m-1}T_k-T_mU_{k-1}=U_{m-k-1}$, since both sides satisfy the same
recurrence in $k$ and agree for $k=0,1$. Take $m=2p$ and use \eqref{eq:cheb-p}:
$T_{2p}=2X^{2p}-1$ and $U_{2p-1}=2T_pU_{p-1}=2X^p(X^2-1)^{(p-1)/2}$. Reflecting at degree $2p+C-1$, the first
term becomes $2XW\operatorname{rev}_{C-1}T_k$, the second $(2-X^{2p})\operatorname{rev}_{C-1}U_{k-1}$, and the
right side $X^{2p+C-1}U_{2p-k-1}(1/X)$. For $k=m_r\ge-g$ this is divisible by $X^{C+k}$, hence by $X^L$, if
$2p-k-1\ge0$. If $2p-k-1<0$ it has degree $k-2p-1$ and is divisible by $X^{4p+C-k}$, hence by $X^L$ when
$k\le4p+g$, which holds because $k\le n-1-g\le4p+g$. Summing with weights $c_r$ and multiplying by
$(X-1)^q(1-X)^h$ gives \eqref{eq:K6}.
\end{proof}

\begin{proof}[Proof of Theorem~\ref{thm:U6}]
Let $M=[\mu_0,\mu_1)$ and $w=\mu_1-\mu_0$. We may assume $w>e$. The block consists of the filtered columns
$k(J,\mu)=Jp+\mu-(b+q)$ for $1\le J\le5$ and $\mu\in M$. These indices are non-negative by the hypothesis
$b+q\le p+\mu_0$, and smaller than $n$ because $5p+\mu_1-1=L-1$. There are $5w$ of them. As $\mu<p$, formula
\eqref{eq:filtered-digits} applies: the response of $c$ at $k(J,\mu)$ is
$\sum_u\alpha_{u,J}Q_c[up+\mu]+\beta_{u,J}\widehat D_c[up+\mu]$.

\emph{Supports.} For $\mu\in M$ we have $p+\mu\le p+\mu_1-1=L-4p-1\le a$ by $n\le4p+2g+1$, so the digits
$u=0,1$ of $Q_c$ and $\widehat D_c$ vanish. From $\mu_0\ge H+q-6p=L+\Lambda-6p$ we get that the digits $u\ge6$ of
$\widehat D_c$ and $u\ge7$ of $Q_c$ vanish. For $u\le5$ we have $up+\mu\le L-1$, so \eqref{eq:K6} applies.

\emph{From $Q$ to $\widehat D$.} For $u=2,\dots,5$, \eqref{eq:K6} gives
$\widehat D_c[up+\mu]=Q_c[up+\mu]+\tfrac12\widehat D_c[(u-2)p+\mu]$, and for $u=2,3$ the last index is at most $a$,
so the term vanishes. Hence if $Q_c[up+\mu]=0$ for $u=2,\dots,5$ and all $\mu\in M$, then by induction
$\widehat D_c[up+\mu]=0$ for $u=2,\dots,5$, and all digits of $\widehat D_c$ in $M$ vanish.

\emph{A relation from the Frobenius map.} We have
$Q_c=(-1)^qX(1-X)^{q+h+(p-1)/2}(1+X)^{(p-1)/2}S_c$, and $q+h+(p-1)/2+e=p$. Since $(1-X)^p=1-X^p$,
\[
 (1-X)^e\,Q_c=(1-X^p)\,Y_c,\qquad Y_c=(-1)^qX(1+X)^{(p-1)/2}S_c,
\]
and $Y_c$ is supported in $[a+1,\,C+(p-1)/2]$. Comparing the coefficients at $up+\mu$, subtracting the
coefficient at $(u-1)p+\mu$, and summing over $u=2,\dots,6$, the right side telescopes to
$Y_c[6p+\mu]-Y_c[p+\mu]$, which vanishes for $\mu\ge\mu_0+e$: the second term because $p+\mu\le a$, the
first because $6p+\mu\ge H+q+e=C+(p+1)/2>\deg Y_c$. So, if $Q_c[up+\mu']=0$ for $u=2,\dots,5$ and all
$\mu'\in M$, then
\[
 \sum_{t=0}^e(-1)^t\binom et\,Q_c[6p+\mu-t]=0\qquad(\mu_0+e\le\mu<\mu_1).
\]
This is a recurrence with unit leading coefficient: if $Q_c[6p+\mu]=0$ for the first $e$ values
$\mu=\mu_0,\dots,\mu_0+e-1$, then it vanishes for all $\mu\in M$.

\emph{Counting.} Impose on $c$ that the $n-5w$ filtered columns outside the block vanish, that
$Q_c[up+\mu]=0$ for $2\le u\le5$ and $\mu\in M$ ($4w$ conditions), and that $Q_c[6p+\mu]=0$ for the first $e$
values of $\mu$. By the three steps above all the responses of the block columns vanish. These are at
most $n-5w+4w+e=n-(w-e)$ linear conditions, so the left kernel has dimension at least $w-e$. Finally,
$w-e$ equals $\mu_1-q-e=n+b+q+h-(11p+1)/2$ if $\mu_0=q$ and $(p-1)/2-g$ if $\mu_0=H+q-6p$.
\end{proof}

\subsection{The second window}
\begin{theorem}\label{thm:S}
Let $e=(p+1)/2-q-h\ge0$ and $n\le4p+2g+1$, and put
\[
 m_0=\max(b+q,\,L+\Lambda-4p),\qquad m_1=\min(p,\,L-3p,\,a+1),
\]
where $a=b+q+2g$. Then $\nu_N(p)\ge m_1-m_0-e$.
\end{theorem}

\begin{proof}
Let $M=[m_0,m_1)$ and $w=m_1-m_0>e$. For $\mu\in M$ we have $\mu\le a$, so the digit $u=0$ of $Q_c$ and
$\widehat D_c$ vanishes. The inequality $\mu\ge L+\Lambda-4p$ gives that the digits $u\ge4$ of $\widehat D_c$
and $u\ge5$ of $Q_c$ vanish. Finally $3p+\mu<L$, so that \eqref{eq:K6} gives, with $D_i=\widehat D_c[ip+\mu]$ and
$Q_i=Q_c[ip+\mu]$,
\[
 Q_1=D_1,\qquad Q_2=D_2,\qquad Q_3=D_3-\tfrac12D_1 .
\]
So every response of the filtered matrix at a column with residue $\mu\in M$ is a combination of the four
functionals $D_1(\mu),D_2(\mu),D_3(\mu),Q_4(\mu)$ of $c$. For $\mu\in M$ the columns $k=\mu-(b+q)$,
$p+\mu-(b+q)$, $2p+\mu-(b+q)$ and $3p+\mu-(b+q)$ exist, since $b+q\le\mu<p$ and $3p+\mu<L$. So there are at least
$4w$ block columns, and their rows lie in the span $F$ of the $4w$ functionals.

We find $w-e$ linear relations among these functionals. As in the proof of Theorem~\ref{thm:U6},
$(1-X)^eQ_c=(1-X^p)Y_c$ with $Y_c$ a polynomial. Put $\sigma(\ell)=\sum_{u\ge0}Q_c[up+\ell]$ and sum the
coefficients of this identity over the residue class of $\mu\ge e$ modulo $p$. The right side gives
$\sum_u\bigl(Y_c[up+\mu]-Y_c[(u-1)p+\mu]\bigr)=0$, since $Y_c$ is a polynomial, and the left side gives
$\sum_{t=0}^e(-1)^t\binom et\sigma(\mu-t)$. For $\mu\in M$ with $\mu-e\ge m_0$ all $\mu-t$ lie in $M$, and there
$\sigma=Q_1+Q_2+Q_3+Q_4=D_1+D_2+D_3-\tfrac12D_1+Q_4$. So
\[
 \sum_{t=0}^e(-1)^t\binom et\bigl(\tfrac12D_1+D_2+D_3+Q_4\bigr)(\mu-t)=0\qquad(m_0+e\le\mu<m_1),
\]
and these $w-e$ relations are independent, since $Q_4(\mu)$ occurs with coefficient $1$ in the one for $\mu$
and only with lower arguments in the others. Hence $\dim F\le4w-(w-e)$, and the rank of the matrix is at most
$(n-4w)+3w+e=n-(w-e)$.
\end{proof}

The proof does not use a determinant or a lower bound for $N$. Unlike the proof of Theorem~\ref{thm:U6}, the
telescoping here is over all digits, so no vanishing of $Y_c$ at the ends is needed.

\subsection{The counting inequality}
The four proofs above have one shape. A space of row combinations is chosen, a few linear conditions are
imposed, and it is shown that the conditions make every response vanish. The dimension of the space minus
the number of conditions is then a lower bound for the nullity. Theorem~\ref{thm:count} isolates this
argument, and allows in addition filters longer than the prime.

\begin{proof}[Proof of Theorem~\ref{thm:count}]
By \eqref{eq:rational-part}, the response of $c$ at the filtered column $k$ is
$\sum_{v=0}^q(-1)^v\binom qvR_c(b+k+v)$, the value at $T$ acting as the shift of the polynomial
$(1-T)^q$. Modulo $p$ we have $(1-T)^p=1-T^p$, hence
\[
 (1-T)^q=(1-T)^r(1-T^p)^a ,
\]
which is the congruence of Lucas's theorem for the binomial coefficients $\binom qv$. So the response at
$k$ equals $\sum_{u=0}^r(-1)^u\binom ru\,S_c(b+k+u)$, where
$S_c(j)=\sum_{v=0}^a(-1)^v\binom av R_c(j+vp)$. The arguments satisfy $b\le j$ and $j+ap\le b+n-1+q=L-1$.
For $c=\iota(x)$ with $x\in\ker\Phi$ we have $S_c(j)=\pm G_x(j)$, and the difference of order $r$ of a
polynomial of degree less than $r$ vanishes identically. So $\iota(\ker\Phi)$ lies in the left kernel of
$\cA^0_{N,p}$, and since $\iota$ is injective, $\nu_N(p)\ge\dim\ker\Phi=\dim\cV-\operatorname{rank}\Phi$.
\end{proof}

If the conditions only force the responses on a set $K$ of filtered columns to vanish, the same argument,
with the remaining $n-|K|$ responses added to the conditions, gives
$\nu_N(p)\ge\dim\cV-\operatorname{rank}\Phi-(n-|K|)$. The proofs of Theorems~\ref{thm:T}--\ref{thm:S} are of
this form, with $K$ the set of all filtered columns in Theorems~\ref{thm:T} and~\ref{thm:U} and a block of
columns in Theorems~\ref{thm:U6} and~\ref{thm:S}.

\begin{proof}[Proof of Theorem~\ref{thm:profile}]
Since $p\nmid\mathcal E$, the reductions of $v_1,\dots,v_d$ are linearly independent, so that $V_p$ is the
reduction of $V$ in the coordinates of this basis. The kernel vectors of $B_j$ satisfy an integer identity,
so their reductions lie in the kernel of $B_j|_{V_p}$, and they remain independent because their minor is a
unit modulo $p$. Hence $\operatorname{rank}(B_j|_{V_p})\le r_j$. The map
$\Phi=(B_j\circ\varphi_{j,k})_{j,k}$ therefore has rank at most $\sum_jw_jr_j$, and its kernel has dimension at
least $md-\sum_jw_jr_j$. Theorem~\ref{thm:count}, applied to $\ker\Phi$ with no further conditions, gives
the bound.
\end{proof}

In the applications the row combinations are organized in \emph{slot systems}. A slot system consists of data
$x_1,\dots,x_D\in\F_p^K$, a linear map $x\mapsto c(x)$ from the data to row combinations, injective on a
subspace of admissible data, fixed polynomials $\lambda_1,\dots,\lambda_D\in\F_p[X]$ of degree less than some
$e_{\max}$, and two linear maps $\mathcal T,\mathcal G$ from $\F_p^K$ to $\F_p[Y]$, such that for every
admissible $x$ the two row polynomials of $c(x)$ multiplied by $(X-1)^h$, respectively by
$(X-1)^h(X^2-1)^{(p-1)/2}$, equal $\sum_d\lambda_d(X)\,\mathcal T(x_d)(X^p)$, respectively
$\sum_d\lambda_d(X)\,\mathcal G(x_d)(X^p)$. This is the form in which \eqref{eq:cheb-p} and the relation
$(1-X)^p=1-X^p$ express the row polynomials modulo $p$. By the digit formula, the response of $c(x)$ at the
raw column $Jp+\ell$ is
\[
 \sum_{0\le e<e_{\max},\ p\mid H-1-\ell-e}\ \psi_{J,\,(H-1-\ell-e)/p}\Bigl(\sum_d\lambda_d[e]\,x_d\Bigr),
\]
where $\lambda_d[e]$ is the coefficient of $X^e$ and each $\psi_{J,s}$ is a linear form, a moment form in $J$
and the digit $s$, determined by $\mathcal T$ and $\mathcal G$. For the stored subspaces the data are $m$
polynomial coordinates with values in $V$, and the admissibility of the data is not imposed directly: the
jets $\varphi_{j,k}$ are Taylor coefficients at $X=-1,0,1$, coefficients at $\infty$, and window rows, the map
$\iota$ from $V_p^m$ to the data is an arbitrary linear map, and only on the joint kernel of the jets is it
required to land in the admissible data, to be injective, and to give responses of the form needed in
Theorem~\ref{thm:count}.

\subsection{Further families}\label{ssec:families}
The following lower bounds for $\nu_N(p)$ are proved, each for every admissible design, by explicit choices of
the space, the conditions and the columns in Theorem~\ref{thm:count} or its variant above. Each is a
piecewise linear function of $p$ and the parameters, positive on windows described by linear inequalities.
We describe them in words. The statements are those of the formalization, cited in
Section~\ref{sec:formal}.
\begin{enumerate}[label=(\roman*),wide]
\item \emph{Frobenius products.} Row combinations whose polynomial is the product of a fixed factor, a
free polynomial $P$ of bounded degree $d_P$ in the Chebyshev variable, and a polynomial $\Lambda(X^p)$ with
exponents in a window of $K$ consecutive values. The conditions require $\Lambda$ to annihilate the
functionals seen by the columns, and the bound is $(d_P+1)(K-\#\text{seen})$.
For $D$ this gives $53N-2p-1$ on $(26N,53N/2)$, where $2p+2\le53N$.
\item \emph{A truncated form of Theorem~\ref{thm:U6}.} Without the hypothesis $n\le4p+2g+1$, the bound
$\mu_1-\mu_0-e_+-(n-4p-2g-1)_+$ holds whenever $L\le6p$ and $b+q\le p+\mu_0$. For $D$ it gives $5p-59N+1$ on
$(59N/5,12N]$.
\item \emph{A raw tail of the antisymmetric family.} With the four centres $g+kp$ ($0\le k\le3$) and
$D_T'=n-1-g-3p$, the count is
$p-1-(2g+q+h)$ under $4p\le L$, $b\le p$, $D_T'\le g$, $D_T'\le(p-1)/2$, $3p+1+g\le n$ and $4p+b\le H$. For
$D$ this is $p-16N-1$ on $(16N,118N/7]$.
\item \emph{Two families of slot polynomials and their sum.} The first fixes a digit $s\le H/p$ and a range of
blocks of rows, with slot polynomials divisible by a power of $X-1$. The second has slot polynomials
$(X^p-1)P$ with $P$ supported in a window of degrees. Their kernels are independent when both apply,
under a separation condition, at the same window $s$ or at the windows $s$ and $s-1$, and the sum of the
two bounds is a lower bound. This is the one place where two bounds are added.
\item \emph{Families with a constant profile on a cell.} Several families (the theorems called MT, M, Z, C
and $\mathrm C'$ in the formalization) use slot polynomials whose coordinates satisfy finitely many
rational linear equations, the \emph{cell} equations. A solution modulo $p$ exists for every $p$ prime to a
fixed integer, the modulus of the cell, and the bound is the dimension of the slot space minus the number of
conditions. In the families C and $\mathrm C'$ the responses depend only on the residue of the column modulo
$p$, so that the filter costs $(p-q)_+$ conditions. In $\mathrm C'$, where $h\ge p$, one writes
$(X-1)^h=(X^p-1)(X-1)^{h-p}$, and the factor $(X-1)^{h-p}$ costs $h-p$ more. The family C gives bounds on $(58N/17,59N/17]$ and
$(29N/8,11N/3]$, and $\mathrm C'$ on $(16N/7,5N/2]$.
\item \emph{Block decompositions for $h\ge p$ and $h<p$.} Writing a row index as $kp+d$ and $y=X^p$, the
identities
\[
 (X-1)^hWT_{kp+d}=\bigl((X-1)T_k\bigr)(y)\,Q_1+\bigl((X-1)(X^2-1)U_{k-1}\bigr)(y)\,Q_2
\]
and its analogue for $U_{kp+d-1}$, with $Q_1=(X-1)^{h-p}WT_d$ and $Q_2=(X-1)^{h-p}U_{d-1}$, hold in $\F_p[X]$ for
$h\ge p$. At a residue with a single live digit the responses of all rows of a block are proportional, and the
bound counts the blocks. For $h<p$ a difference of order $p-h$ of the coefficients plays the same role.
For $D$ these give bounds on $(59N/21,3N]$, where $h\ge p$, and on $(3N,16N/5]$.
\end{enumerate}

\subsection{Stored data for finitely many windows}\label{ssec:packets}
On the remaining windows below $16.2N$ the bound is that of Theorem~\ref{thm:profile} with a stored
subspace. For each of $34$ windows we store integer vectors $v_1,\dots,v_d$ spanning $V$, a set of $d$ pivot
columns with the adjugate of the corresponding minor, and for each map $B_j$ its integer matrix, the bound
$r_j$, $d-r_j$ kernel vectors in the coordinates of $V$, and again a pivot set with an adjugate. A single
Boolean function of these data, evaluated by the Lean kernel, checks the shapes, the identities
$\mathrm{adj}\cdot v_{\mathrm{piv}}=\delta\cdot I$ and the products $B_j\bigl(\sum_ak_av_a\bigr)=0$ in
integers. Its soundness theorem gives the hypotheses of Theorem~\ref{thm:profile} on $V$ and the $B_j$, with
$\mathcal E$ the product of the minors. No prime enters this check.

The hypotheses on the jets and on $\iota$ depend on $N$ and $p$. A window is described by finitely many
inequalities that are affine in $N$ and $(p-1)/2$, and so are $m$ and the $w_j$ on each of finitely many
pieces of the window. A second verified procedure checks on each piece that the numbers of jets and their
orders at $-1$, $0$, $1$ and $\infty$ are nonnegative and fit the window, that $p>59$, that the map $\iota$
is injective on the joint kernel and lands in the admissible data of a slot system, and that the response is
given on the joint kernel by a polynomial in the column of degree less than $q \bmod p$ (so that
Theorem~\ref{thm:count} applies), each required inequality in $N$ and $(p-1)/2$ being derived from the
window inequalities by sums, nonnegative integer multiples, and division of the variable part by a common
divisor with rounding of the constant, that is, by Chv\'atal--Gomory cuts. Its soundness theorem gives the
hypotheses of Theorem~\ref{thm:profile} at every odd prime $p>59$ of the window with $p\nmid\mathcal E$.
For large $N$ all primes of the window are of this kind. For $33$ of the $34$ windows the hypotheses are
established in this way, and for the remaining one by a direct proof of the same statement. The windows
lie in $(3N,16.2N)$, where $p>q$, so $a=0$ in Theorem~\ref{thm:count}.

\subsection{The function $\kappa$ for the design $D$}\label{ssec:kappa}
For $D$, $L=70N$, $H=78N$, $a=22N$, $\Lambda=11N$ and $q+h=6N$. On $(12N,53N/2]$ the function $\kappa$ of
Theorem~\ref{thm:profile-D} is as follows, with $x=p/N$.
\begin{center}
\begin{tabular}{@{}lll@{}}
\toprule
interval of $x$ & $\kappa(x)$ & bound \\
\midrule
$(12,13]$ & $x/2-5$ & Theorem~\ref{thm:content}(c), $(p-1)/2-g$ \\
$(13,\,146/11]$ & $73-11x/2$ & Theorem~\ref{thm:content}(c), second term \\
$(341/25,\,69/5]$ & $11x/2-75$ & Section~\ref{ssec:families}(v), family Z \\
$(69/5,\,14]$ & $x/2-6$ & Section~\ref{ssec:families}(v), family Z \\
$(14,\,711/50]$ & $64-9x/2$ & Section~\ref{ssec:families}(v), family Z \\
$(30287/2000,\,78/5]$ & $7x/2-53$ & Section~\ref{ssec:families}(v), family Z \\
$(78/5,\,16]$ & $25-3x/2$ & Section~\ref{ssec:families}(v), family Z \\
$(16,\,809/50]$ & $73-9x/2$ & Section~\ref{ssec:families}(v), family Z \\
$(809/50,\,178/11]$ & $73-9x/2$ & Section~\ref{ssec:packets} \\
$(178/11,\,1619/100]$ & $x-16$ & Section~\ref{ssec:packets} \\
$(1619/100,\,118/7]$ & $x-16$ & Section~\ref{ssec:families}(iii) \\
$(118/7,\,69/4]$ & $9x/2-75$ & Theorem~\ref{thm:content}(d), $(9p-2C)/2$ \\
$(69/4,\,35/2]$ & $x/2-6$ & Theorem~\ref{thm:content}(d), $(p-1)/2+h-b$ \\
$(35/2,\,128/7]$ & $64-7x/2$ & Theorem~\ref{thm:content}(d), $(2(n+q+h)-7p)/2$ \\
$(64/3,\,23]$ & $3x-64$ & Theorem~\ref{thm:content}(a), $3p-(n+q+h)$ \\
$(23,\,24]$ & $5$ & Theorem~\ref{thm:content}(a) up to $x=70/3$, then (b) \\
$(24,\,26]$ & $53-2x$ & Theorem~\ref{thm:content}(b), $n-g-2p$ \\
$(26,\,53/2]$ & $53-2x$ & Section~\ref{ssec:families}(i) \\
\bottomrule
\end{tabular}
\end{center}
The six rows on $(341/25,809/50]$ come from the family Z of Section~\ref{ssec:families}(v), in two forms, on
$(341/25,711/50]$ and on $(30287/2000,809/50]$. The areas of the rows of the table add up to $21.02610\ldots$,
and the $84$ intervals below $12$ contribute $1.03361\ldots$, so that
the total is $\Chi\ldots$. The exact sum is computed in the formalization from the list of intervals.

For Theorem~\ref{thm:content} one checks the hypotheses on the intervals above. The hypotheses
$e\ge0$ and $n\le4p+2g+1$ both read $x\ge12$, and $L\le6p$ reads $x\ge35/3$. So Theorems~\ref{thm:U6}
and~\ref{thm:S} apply for $x\ge12$. For $12\le x\le13$ one has $\mu_0=H+q-6p=81N-6p$, so the first branch
of Theorem~\ref{thm:U6} applies, and for $x\ge13$ one has $\mu_0=3N$ and the second branch, which is positive
until $x=146/11$. In Theorem~\ref{thm:S}, $m_1=p$ for $x\le35/2$ and $m_1=70N-3p$ above, and
$m_0=81N-4p$ for $x\le69/4$ and $12N$ above. The value $m_1-m_0-e$ is $9x/2-75$, $x/2-6$ and
$64-7x/2$ on the three intervals $(50/3,69/4]$, $(69/4,35/2]$, $(35/2,128/7]$, and it is non-positive
outside $(50/3,128/7)$. Theorem~\ref{thm:T} applies for $x>10$ and $x\ge39/2$, and
$D_T=\min(5N,52N-1-2p,75N-3p)$. It gives $3p-(n+q+h)=3x-64$ up to $x=23$ and the cap $g=5N$ up to
$x=70/3$. From $x=70/3$ the hypotheses of Theorem~\ref{thm:U} hold up to $x=26$ (they read
$3p\ge L$, $3p\le H$, $2p\le n+q$ and $2h\le p+1$), and give $\min(n,2p+2g+1)-g-2p-1$, which is $5N$ for
$x\le24$ and $53N-2p-1$ above. In each case the value of the theorem is the value in the table minus at
most one unit, from the integer rounding. The same kind of check, with linear inequalities in $N$, $p$
and the parameters, gives the bounds of Section~\ref{ssec:families} on their intervals, and
Section~\ref{ssec:packets} gives the others. This proves Theorem~\ref{thm:profile-D}. The loss of $2$ in its
statement covers the integer rounding of all the bounds.

\begin{remark}
Numerical computations of the exact rank of $\cA^0_{N,p}$ for several designs and primes show that the
bounds of Theorem~\ref{thm:content} are attained in the windows where they were tested, and that other
families add kernel where the windows overlap. These computations are not part of the proof. An upper bound
for the nullity is not proved, and none is needed.
\end{remark}

\begin{remark}
The subspaces $V$ of Section~\ref{ssec:packets} were found numerically, as optimal constant subspaces for the
bound of Theorem~\ref{thm:profile} on each window. They are dense rational subspaces with no evident closed
form, which is why they are stored. Only the stored integer identities and the checked window conditions
enter the proof.
\end{remark}
\section{The real place: perturbation and fixed columns}\label{sec:real}

In this section and the next we prove Theorem~\ref{thm:real}. The determinant $\Delta_N$ must be small
not only at $G$ but on an interval of length $e^{-\rho N}$ around it. Here we reduce this to determinants
in which some columns of $\cA_N$ are replaced by values of the row functions at fixed points
(\S\ref{ssec:quadrature}--\ref{ssec:andreief}). We then recall OpenAI's energy method, extended to
configurations with fewer free points (\S\ref{ssec:energy}). It splits into two cases. In the first case it
already gives the decay rate of Theorem~\ref{thm:real}. In the second case it falls short by the gain
$\Gain$ when few columns are fixed. Section~\ref{sec:signed} closes that gap. The section ends with the
trial sequences and the verification of the bounds (\S\ref{ssec:trials}--\ref{ssec:rho}).

The design is a parameter throughout. We write $\alpha_A=A/n=(b_0+q_0+2g_0)/n_0$, $\beta_h=2h_0/n_0$,
$\gamma_H=H/n=(n_0+b_0+q_0+g_0+h_0)/n_0$, $\gamma_d=d_0/n_0=(n_0-g_0)/n_0$, and
\[
 c_{07}=1-\frac{b_0+q_0+g_0}{n_0},\qquad c_{26}=3-\frac{b_0+q_0+2g_0}{n_0}.
\]
For $D$ these are $\alpha_A=11/29$, $\beta_h=3/29$, $\gamma_H=39/29$, $\gamma_d=53/58$, $c_{07}=41/58$ and
$c_{26}=76/29$.

\subsection{A quadrature formula for the coefficient of the constant}\label{ssec:quadrature}
Let $\delta=x-G$. Then $\Delta_N(x)=\det(\cA_N+\delta\cL_N)$. The question is how large the columns of
$\cL_N$ are compared with those of $\cA_N$. By \eqref{eq:L} the entries of $\cL_N$ are linear functionals of
the coefficients of the row polynomials. These polynomials have bounded degree. So the functionals can be
written exactly as finite combinations of point evaluations. If $\tau_1,\ldots,\tau_J$ are distinct points
and $J$ is at least the number of degrees of freedom of the rows, Lagrange interpolation gives
\begin{equation}\label{eq:quad}
 \cL_N(r,k)=\sum_{a=1}^JA_r(\tau_a)\,\beta_{a,k},
\end{equation}
where the weight $\beta_{a,k}$ is the value of the same functional on the $a$-th Lagrange basis polynomial.
This is a Gauss-type quadrature formula. It makes the columns of $\cL_N$ combinations of point evaluations of
the row functions that define $\cA_N$. The perturbed determinant can then be treated by the energy method.
The cost is the size of the weights.

We take $J=d_0N=53N$ nodes. Put $g_*=401/2500$ and $a_*=2g_*/(1+g_*^2)=2005000/6410801$, and let
$K=\lceil J/2\rceil$. The nodes are the points
\[
 t_{j,\pm}=\pm\sqrt{a_*^2+(1-a_*^2)(1+s_j)/2},\qquad T_K(s_j)=0,
\]
that is, the roots of $T_K\bigl((2t^2-(1+a_*^2))/(1-a_*^2)\bigr)$. We keep all $K$ positive roots and the
first $J-K$ negative ones, so that for odd $J$ exactly one root is dropped. The nodes lie in
$S=(-1,-a_*)\cup(a_*,1)$. In the coordinate $x$ with $t=2x/(1+x^2)$ every node has $|x|>g_*$. This matters
below, because the energy fields are larger near $x=0$. Dropping a root changes the nodal polynomial by a
linear factor. It changes the derivative bound and the size of the Lagrange polynomials by factors
polynomial in $N$ only. So an odd number of nodes costs nothing in the exponent. This is needed, since
$J=53N$ is odd for odd $N$.

Let $g_S$ be the Green function of $\C\setminus\bar S$ with pole at infinity. In closed form
\[
 g_S(z)=\tfrac12\operatorname{arccosh}\Bigl(\frac{|z^2-1|+|z^2-a_*^2|}{1-a_*^2}\Bigr).
\]
The Lagrange polynomials satisfy $|\ell_a(z)|\le e^{53Ng_S(z)+o(N)}$ at distance at least $c/N$ from the
nodes. This follows from the derivative identity for Chebyshev polynomials at their roots. The weight of
column $k$ is a coefficient functional of $\ell_a$ multiplied by the column filter. After removal of the
common factor of the rows it equals
\[
 \lambda_k=4\,[t^0]\,F_k(t)\,(1-t^{-2})^{-1/2},\qquad F_k(t)=t^{\,A-b-q-k}(1-t)^h(t-1)^q\ell_a(t),
\]
where $(1-t^{-2})^{-1/2}=\sum_lc_lt^{-2l}$ with $c_l=\binom{2l}l4^{-l}$. The next lemma bounds it by a
contour that depends on $k$.

\begin{lemma}[keyhole bound]\label{lem:keyhole}
Let $P$ be a real polynomial, $J\ge0$, and $F(t)=t^{-J}P(t)$. Let $0<r\le1$, and suppose that $|F|\le M$ on
the circle $|t|=r$ and $|F|\le M'$ on the segments $r\le|x|\le1$. Then
\[
 \Bigl|4\sum_{l\ge0}c_l\,[t^{J+2l}]P\Bigr|\le4(1+J)\max(M,M').
\]
\end{lemma}

\begin{proof}
Write $F=\Pi+R$ with $\Pi$ a polynomial and $R=\sum_{i=1}^J\pi_it^{-i}$ the principal part at $0$. The
functional sees only $\Pi$. Since $c_l=\frac1\pi\int_{-1}^1x^{2l}(1-x^2)^{-1/2}dx$, its value is
$\frac4\pi\int_{-1}^1\Pi(x)(1-x^2)^{-1/2}dx$, which is at most $4\max_{[-1,1]}|\Pi|$. By Cauchy's estimate
$|\pi_i|\le Mr^i$, so $|R(x)|\le JM$ for $r\le|x|\le1$, and $|\Pi|\le M'+JM$ there. For $|x|<r$ the maximum
principle gives $|\Pi(x)|\le\max_{|t|=r}|\Pi|\le(1+J)M$.
\end{proof}

In the formal proof the integral is replaced by the finite Gauss--Chebyshev identity
$\frac1M\sum_{j<M}\cos^{2l}\frac{(2j+1)\pi}{2M}=c_l$ for $l<M$, and the Cauchy estimate by the finite Cauchy
formula at roots of unity. So no integration is needed. With $\kappa=k/n$, Lemma~\ref{lem:keyhole} gives
\begin{equation}\label{eq:wk}
\begin{gathered}
 |\lambda_k|\le e^{(w(\kappa)+o(1))N},\\
 w(\kappa)=\max_{z\in\Gamma_r}\bigl[(10-58\kappa)\log|z|+6\log|1-z|+53g_S(z)\bigr],
\end{gathered}
\end{equation}
where $\Gamma_r$ is the circle $|z|=r$ together with the segments $r\le|x|\le1$. Here $10=(A-b-q)/N$ and
$6=(q+h)/N$. The radius $r$ may depend on $\kappa$. For $\kappa$ near $1$ the factor $|z|^{10-58\kappa}$
favours a large circle, and for small $\kappa$ a small one. A single circle of radius $1.001$ gives
$w(\kappa)\le52.785$ for every $\kappa$. The keyhole contour gains only a little over it at the worst
column. Its value is that the worst column for the weight is not the worst column for the binomial factor
below.

\subsection{Column expansion and an identity of Andr\'eief type}\label{ssec:andreief}
Write $\Phi(r,a)=A_r(\tau_a)$. Expanding $\det(\cA_N+\delta\Phi\beta)$ multilinearly in the columns gives a
sum over maps $f$ from the columns to $\{\varnothing,1,\ldots,J\}$ of
$\prod_{f(k)\ne\varnothing}\delta\beta_{f(k),k}\det M_f$. Here $M_f$ is the matrix whose $k$-th column is the
$k$-th column of $\cA_N$ if $f(k)=\varnothing$, and the vector $(A_r(\tau_{f(k)}))_r$ otherwise.

\begin{lemma}\label{lem:expansion}
Suppose that $|\det M_f|\le M\prod_{f(k)\ne\varnothing}C_k$ for every $f$, and that $|\beta_{a,k}|\le B_k$. Then
\[
 \bigl|\det(\cA_N+\delta\Phi\beta)-\det\cA_N\bigr|\le M\Bigl(\prod_k(1+J|\delta|B_kC_k)-1\Bigr).
\]
\end{lemma}

The lemma allows a different constant for each column. Each $M_f$ has $p=n-m$ columns that are integrals and
$m$ columns that are point values. Andr\'eief's identity extends to this situation.

\begin{lemma}\label{lem:andreief}
Let $\mu,\nu$ be finite measures, $F_r$ $(r<p+m)$, $G_i$ $(i<p)$ and $K$ integrable functions, and let $B$
be the matrix with $B(r,i)=\iint F_r(t)K(t,s)G_i(s)\,d\mu(t)\,d\nu(s)$ for $i<p$ and
$B(r,p+j)=F_r(\tau'_j)$. Then
\begin{multline*}
 (p!)^2\det B=\int\!\!\int\det\bigl[F_r((t\Vert\tau')_c)\bigr]_{r,c}\\
 \cdot\det[K(t_i,s_j)]_{i,j<p}\det[G_i(s_j)]_{j,i<p}\,d\mu^p(t)\,d\nu^p(s),
\end{multline*}
where $t\Vert\tau'$ is the list of the $p$ variables followed by the $m$ fixed points.
\end{lemma}

\begin{proof}
For $m=0$ this is Andr\'eief's identity applied twice \cite{Forrester2018Andreief}. In general, expand
$\det B$ multilinearly in its first $p$ columns. This introduces one pair of variables $(t_i,s_i)$ for each
of them. Then antisymmetrize, first in the $t$ variables and then in the $s$ variables. The last $m$ columns
are constant throughout.
\end{proof}

With $K(t,s)=1/(1-ts)$ the middle determinant is a Cauchy determinant. The third is a generalized
Vandermonde determinant $\prod_js_j^b(1-s_j)^q\det[s_j^{E_i}]$, where the exponents $E_i<n$ are those of the
columns that remain integrals. The missing exponents are the replaced columns. We bound the generalized
Vandermonde determinant by a product over the missing exponents.

\begin{lemma}\label{lem:binomvand}
Let $s_1,\ldots,s_p\in[0,1]$, let $E_1<\dots<E_p<n$ be integers, and let $\{k_1,\ldots,k_m\}$ be the
complement of $\{E_i\}$ in $\{0,\ldots,n-1\}$. Then
\[
 \bigl|\det[s_l^{E_i}]\bigr|\le\prod_{j=1}^m\beta(k_j)\prod_{i<l}|s_i-s_l|,\qquad
 \beta(k)=\sum_{i=0}^{m-1}\binom p{p-k+i},
\]
with $\binom pr=0$ for $r\notin[0,p]$.
\end{lemma}

\begin{proof}
Let $P(X)=\prod_l(X-s_l)=\sum_r(-1)^re_r(s)X^{p-r}$. Multiplying the matrix by a unitriangular matrix built
from $P$ gives $|\det[s_l^{E_i}]|=\prod_{i<l}|s_i-s_l|\cdot|\det Q|$, where $Q$ is the $m\times m$ matrix with
entries $Q_{j,i}=[X^{k_j-i}]P$. So $|Q_{j,i}|=e_{p-k_j+i}(s)\le\binom p{p-k_j+i}$, since $s\in[0,1]^p$.
Hadamard's inequality, applied to the rows of $Q$ in the $\ell^1$ norm, gives the claim.
\end{proof}

The bound does not depend on the points $s_l$. Since $\binom pr\le\binom nr$, one has
$\log\beta(k)\le n_0\tilde H(k/n)N+O(\log N)$, where $\tilde H(\kappa)$ is the binary entropy $H(\kappa)$ for
$\kappa\ge1/2$ and $\log2$ for $\kappa<1/2$. The weight \eqref{eq:wk} grows with $\kappa$ up to
$\kappa\approx0.84$, while $\tilde H$ is largest at $\kappa=1/2$. So the product of the two factors is
smaller than the product of their maxima.

\subsection{The energy estimate with fixed points}\label{ssec:energy}
OpenAI bounds the integrand in Andr\'eief's identity by the exponential of $n^2E$, where $E$ is a discrete
logarithmic energy of the points $x_i$ and $s_j$ in the coordinate $t=2x/(1+x^2)$. The energy consists of an
external field for each point, logarithmic interactions $\log|x_i-x_j|$, $\log(1-x_ix_j)$ and
$\log|s_i-s_j|$ among points of the same kind, and an interaction $-\log(1-2x_is_j+x_i^2)$ between the two
kinds. There are two cases, according to which of OpenAI's two bounds for the determinant of the row
functions is used \cite[Props.~5.2, 5.3]{OpenAI2026}. In the second case the interaction among the $x_i$ has
weight $2$ instead of $1$. In the first case the row field acquires a term with a multiplier $\lambda_1$ for
the constraint that defines the case. The second case also bounds a determinant $\det\mathsf L$ of size $n$
by $6^n$. This is the factor that Section~\ref{sec:signed} improves.

The energy is bounded from above by duality. Given trial sequences $u=(u_k)$ and $v=(v_k)$, one writes the
interaction terms as quadratic forms in the moments $\frac1n\sum_iT_k(x_i)$ and $\frac1n\sum_js_j^k$. Then one
completes the squares against $u$ and $v$. What remains is the sum of the norms $\|u\|^2=\sum_ku_k^2/k$ and
$\|v\|^2$ and, for each point, a function of one variable. With $\mathsf T(u,x)=\sum_ku_kT_k(x)/k$ and
$\mathsf S(v,x)=\sum_kv_kx^k/k$ these are, in the second case,
\begin{align*}
 X(x)&=\alpha_A\log|x|+\beta_h\log(1-x)-\gamma_H\log(1+x^2)-4\mathsf T(u,x)-\mathsf S(v,x),\\
 Y(s)&=\tfrac{b_0}{n_0}\log s+\tfrac{q_0}{n_0}\log(1-s)+2\mathsf T(v,s).
\end{align*}
In the first case they are
\begin{multline*}
 X_1(x)=\alpha_A\log|x|+\beta_h\log(1-x)-(\gamma_H-\tfrac12)\log(1+x^2)\\
 +\lambda_1\bigl(\tfrac{2g_0}{n_0}-\tfrac{2x^2}{1+x^2}\bigr)-2\mathsf T(p,x)-\mathsf S(v,x),
\end{multline*}
with the same $Y$. The trial sequence for the first variables is called $p$ in this case. If $X\le M_X$ and
$Y\le M_Y$ on the whole domains, the energy of every configuration is at most a constant determined by $M_X$,
$M_Y$ and the norms. In potential-theoretic terms the extremal configuration is described by an equilibrium
measure for two interacting measures in an external field \cite{SaffTotik}. The best trial sequences are its
moments, as Carreira observed \cite{Carreira2026}. Any trial sequences give a valid upper bound.

When $m$ of the $n$ columns are point values, the configuration has only $n-m$ free points of each kind. The
normalization is still by $n$. The fixed points enter through the determinant of the row functions evaluated
at the list $t\Vert\tau'$. They are of two kinds. Some are quadrature nodes, with $|x|>g_*$. The others are
points $y\in(0,1)$. We write $\alpha=m/n$ for the fraction of fixed columns and $\tau$ for the fraction of
fixed points of the second kind. The missing free points change the energy bound by an amount linear in
$m$. This is the energy cost of a fixed column. Every step of OpenAI's argument remains valid uniformly in
the number of free points. The regularization of the logarithmic interactions, the diagonal corrections and
the completion of squares only involve sums over the points present. For the fixed points of the first kind
the relevant field is the residual field
\[
 R(x)=-\gamma_d\log(1+x^2)-2\mathsf T(u,x),\qquad g_*\le|x|\le1 .
\]
For those of the second kind it is
\[
 N(y)=-\gamma_d\log(1+y^2)+\frac{n_0-2g_0-1}{n_0}\log y-u_1y,\qquad0<y<1,
\]
where $u_1$ is the first coefficient of $u$. The coefficient $(n_0-2g_0-1)/n_0$ is a lower bound, valid for
all $N\ge1$, of $(n-1-2g)/n$.

With Lemma~\ref{lem:binomvand} the generalized Vandermonde determinant contributes only the factor
$\prod_{i<l}|s_i-s_l|$ to the energy. The binomial factors are charged to the columns. The energy bound then
reduces to a quadratic inequality in $\alpha$ and $\tau$.

\begin{lemma}[budget, second case]\label{lem:budget2}
Let $\ell=\log2$, and let $\mathcal N,\Phi,\mathcal R,\mathcal N_c$, $c_2$ and $m\ge0$ be real numbers with
\begin{gather*}
 \mathcal N+\Phi+U-c_{07}\ell\le-m,\qquad \ell-\mathcal R+\mathcal N_c\le0,\\
 \bigl((c_{26}-1)\ell-\Phi+\mathcal R-c_2\bigr)^2\le4\ell m .
\end{gather*}
Then for all $\alpha,\tau\ge0$
\begin{multline*}
 \bigl(-(c_{07}-\tfrac12)+(c_{26}-1)\alpha-\alpha^2+\tau-\tau^2\bigr)\ell+\mathcal N+(1-\alpha)\Phi\\
 +(\alpha-\tau)\mathcal R+\tau\mathcal N_c\le-U+\tfrac12\ell+c_2\alpha .
\end{multline*}
\end{lemma}

\begin{proof}
The left side minus the right side equals
$(\mathcal N+\Phi+U-c_{07}\ell)+\alpha\bigl((c_{26}-1)\ell-\Phi+\mathcal R-c_2\bigr)-\ell\alpha^2
+\tau(\ell-\mathcal R+\mathcal N_c)-\ell\tau^2$. The last two terms are non-positive by the third hypothesis.
The first three form a concave quadratic in $\alpha$. It is non-positive for all $\alpha$, because
$a+s\alpha-\ell\alpha^2\le0$ for all $\alpha$ as soon as $s^2+4\ell a\le0$.
\end{proof}

The lemma is applied with $\mathcal N$ a bound for $2\|u\|^2+\frac12\|v\|^2$, $\Phi=M_X+M_Y$, $\mathcal R$ a
bound for the supremum of $R$ and $\mathcal N_c$ a bound for the supremum of $N$. The analogue for the first
case has no fixed points of the second kind. Its polynomial in $\alpha$ is
\[
 \bigl(-(c_{07}-\tfrac12)+(c_{26}-1)\alpha-\alpha^2\bigr)\ell+n_1+(1-\alpha)F_1+\alpha\lambda_1\tfrac{2g_0}{n_0}
 \le-U+\tfrac12\ell+c_1\alpha,
\]
with $n_1$ a bound for the norms of the first trial and $F_1=M_{X_1}+M_Y$. It holds for all real $\alpha$
when its discriminant is non-positive and $c_1\le(c_{26}-1)\ell-F_1+\lambda_1\tfrac{2g_0}{n_0}$.

\subsection{The two cases at the improved rate}\label{ssec:cases}
Theorem~\ref{thm:real} asks for the rate $U_0=U+\Gain/n_0^2$ instead of $U$. A larger target is a smaller
slack $m$ in Lemma~\ref{lem:budget2}. In the first case the trial of \S\ref{ssec:trials} has enough slack:
its budget holds at $U_0$ for every $\alpha$, with the cost $c_2$ of the second case. In the second case the
budget at $U_0$ fails for small $\alpha$ and holds for $\alpha\ge1/16$. The reason is that the constant
$U_0-U$ is absorbed by the linear term once $\alpha$ is not small. So only the second case with fewer than
$n/16$ fixed columns needs a new argument. That argument is Section~\ref{sec:signed}. It proves the bound of
the second case at $U_0$ with the same cost $c_2$, for $0\le\alpha\le1/16$.

Assume this for the moment. A determinant $M_f$ with $m$ fixed columns $k_1,\ldots,k_m$ is then at most
$C^m\prod_j\beta(k_j)e^{n_0c_2Nm}$ times $\exp\bigl(n_0^2N^2(-U_0+\frac12\log2+\varepsilon)\bigr)$, with an
absolute constant $C$. The factor $C^m$ comes from the expansion of the row functions at the fixed points.
By \eqref{eq:wk} and the bound for $\beta(k)$, the constants of Lemma~\ref{lem:expansion} satisfy
\[
 \log(B_kC_k)\le\bigl(n_0c_2+w(k/n)+n_0\tilde H(k/n)+o(1)\bigr)N\le(\rho+o(1))N,
\]
with $\rho=n_0c_2+\Psi_0$ and $\Psi_0\ge\sup_\kappa\bigl(w(\kappa)+n_0\tilde H(\kappa)\bigr)$. For
$|\delta|\le e^{-\rho N}$ Lemma~\ref{lem:expansion} gives a factor $(1+O(N^5))^n=e^{o(N^2)}$. This proves
Theorem~\ref{thm:real}, given Section~\ref{sec:signed} and the numerical bounds below.

\subsection{The trial sequences}\label{ssec:trials}
For the first case we use a trial sequence with $40$ Chebyshev coefficients and a finite number of geometric
terms, with the multiplier $\lambda_1=1.68493357$. For the second case we use trial sequences of the form
\begin{equation}\label{eq:trial}
 u_k=c_k+\operatorname{Re}\sum_jr_jz_j^k\qquad(k\ge1;\ c_k=0\text{ for }k>K),
\end{equation}
with $K=\Kcoef$ Chebyshev coefficients, which are dyadic rationals, and $\Ngeo$ geometric terms. Of these,
$\Ncomplex$ have complex $z_j$ and complex $r_j$.

The geometric terms have a simple meaning. Since $\sum_{k\ge1}z^kT_k(x)/k=-\frac12\log(1-2xz+z^2)$, a term
$(z,r)$ contributes $-\frac12\operatorname{Re}\,r\log(1-2xz+z^2)$ to the potential $\mathsf T(u,x)$. If
$x_0+i\epsilon=\frac12(z+z^{-1})$, then $1-2xz+z^2=2z(x_0+i\epsilon-x)$. So the term is, up to a constant, the
logarithmic potential of a point charge of complex strength $r$ at $x_0+i\epsilon$ and its conjugate. The
extremal measure has a gap in its support around $x=0$. Its potential has a logarithmic singularity at $x=0$
and square-root behaviour at the edges of the gap. A Chebyshev expansion converges slowly near such points.
A charge at a small distance $\epsilon$ reproduces the local behaviour with one term. We place the complex
charges at the $41$ equally spaced points $x_0=-0.16+0.008j$ $(0\le j\le40)$, all with $\epsilon=10^{-2}$.
The real terms sit at points $\pm z$ with $z$ increasing from $0.8$ to $0.999997$. The numbers $z_j$ are
rationals with at most twelve decimals. The norms of the geometric terms have closed forms. Two terms
$(z_j,r_j)$ and $(z_l,r_l)$ contribute
$-\frac12\operatorname{Re}\bigl(r_jr_l\log(1-z_jz_l)+r_j\bar r_l\log(1-z_j\bar z_l)\bigr)$ to $\|u\|^2$.

\subsection{Verification of the bounds}\label{ssec:cells}
We must bound three functions of one variable from above on their whole domains. They are the field $X$ on
$[-1,1]$, the field $Y$ on $[0,1]$, and the residual field $R$ on $g_*\le|x|\le1$. With $x=\cos\theta$, each is
a constant plus a cosine series with $K$ terms plus finitely many terms $\operatorname{Re}\,w\log(1-ae^{im\theta})$.

We cover each domain by cells $[x_c-h,x_c+h]$ whose endpoints are multiples of $2^{-40}$. On each cell we
bound the function by its Taylor expansion at $x_c$. The cells are grouped in blocks. For each block the
derivatives of the polynomial part are computed exactly at the centre of the block up to order $70$. The
remainder beyond order $70$ is bounded by a single constant for the block, multiplied by $h^{70}/70!$. The
Taylor expansion at the centre of a cell is obtained from that of the block by an exact shift. The
logarithmic terms are expanded to order $5$, with remainders bounded from the exact minimum of
$|1-ae^{im\theta}|$ on the cell. A few cells at $x=0$ and $x=\pm1$ treat the singular terms separately. The
covers of $X$, $Y$ and $R$ have \Nblocks{} blocks and \Ncells{} cells. Each cell is checked by exact rational and
integer arithmetic.

The resulting bounds are
\[
 \sup X\le\SupX,\quad \sup Y\le\SupY,\quad \sup_{g_*\le|x|\le1}R\le\ResCap .
\]
So $M_X+M_Y=\FieldCap$, and the norms satisfy $2\|u\|^2+\frac12\|v\|^2\le\NormCap$. The first coefficient
satisfies $u_1\ge9/50$, which is used near $x=0$. For the points of the second kind a one-variable estimate
shows $\sup_{0<y<1}N(y)\le-77/100$ for $u_1\ge9/50$. For the first case, with the weights
$(\alpha_A,\beta_h,\gamma_H-\frac12,\lambda_1,2g_0/n_0)=(11/29,3/29,49/58,\lambda_1,5/29)$, the bounds
$\sup X_1\le-1.15854$ and $\sup Y\le-1.32981$ are obtained by isolating the zeros of the numerators of the
derivatives. These numerators are polynomials with integer coefficients. The norm bound is $n_1=0.67761$.

With $c_{07}=41/58$, $c_{26}=76/29$ and $U=\Utarget$, the hypotheses of Lemma~\ref{lem:budget2} hold with
$m=\Mmargin$ and $c_2=\Ctwo$. The first case admits $c_1=\Cone<c_2$. At the target $U_0$ the first-case
budget still holds with the cost $c_2$, and the second-case budget holds for $\alpha\ge1/16$ with the cost
$c_2$. The trial data are fixed once and for all. The file containing them has SHA-256 digest
\begin{center}
{\footnotesize\texttt{\Digest}}.
\end{center}

\subsection{The perturbation exponent}\label{ssec:rho}
It remains to bound $\Psi_0$. For a fixed radius the function $\kappa\mapsto\max_{\Gamma_r}\Phi_\kappa$ is a
maximum of functions affine in $\kappa$, hence convex. The entropy $H$ is concave. On a cell
$[\kappa_1,\kappa_2]\subset[\frac12,1]$ we replace $H$ by its tangent at the midpoint. The sum is then
convex, and its supremum on the cell is attained at an endpoint. On $[0,\frac12]$ the term $\tilde H=\log2$
is constant. We cover $[0,1]$ by $14$ cells, each with a rational radius $r$. At each endpoint the maximum of
$\Phi_\kappa$ on the circle and on the segments is bounded by a cover of the contour by small arcs and
intervals, in exact arithmetic. This proves
\[
 \sup_{0\le\kappa\le1}\bigl(w(\kappa)+58\tilde H(\kappa)\bigr)\le\Psi_0=\PsiZero .
\]
The supremum is attained near $\kappa=0.62$, with radius about $0.6$. Hence $\rho=58\cdot\Ctwo+\PsiZero=\Rho$.
\section{The real place: a signed bound for the second case}\label{sec:signed}

In this section we prove the bound of the second case at the rate $U_0$, for configurations with at most
$n/16$ fixed columns. This completes the proof of Theorem~\ref{thm:real}. The section is organized as
follows. The second case contains a determinant $\det\mathsf L$ of size $n$ that the energy method bounds by
$6^n$. We first write it with Pick matrices (\S\ref{ssec:pick}). A Schur complement argument and a contour
bound reduce it to the block of the negative nodes (\S\ref{ssec:schur}--\ref{ssec:T}). Heine's formula turns
that block into a ratio of two integrals of log-gas type (\S\ref{ssec:heine}). We bound the numerator from
above and the denominator from below (\S\ref{ssec:loggas}). Positive nodes close to $1$ need a separate
argument (\S\ref{ssec:spikes}). Combined with OpenAI's energy, the bounds give a concave majorant thanks to an
inequality for signed logarithmic energy (\S\ref{ssec:reflection}). Fixed columns add fixed rows
(\S\ref{ssec:fixedrows}). The hypotheses of the reduction hold on all configurations that matter
(\S\ref{ssec:class}). The section ends with the trials and their verification (\S\ref{ssec:signedcert}).

\subsection{The determinant of the row functions}\label{ssec:pick}
In the second case OpenAI's integrand contains the determinant of the row functions at the points
$x_1,\ldots,x_n\in(-1,1)\setminus\{0\}$ \cite[Prop.~5.3]{OpenAI2026}. With $e=n-1-2g$ it factors as
\begin{equation}\label{eq:detphi}
 \det\varphi(x)=c\,V(x)\prod_a|x_a|^{-(n-1-g)}\det\mathsf L(x),
\end{equation}
where $c$ depends only on $n$ and $g$, and $V(x)=\prod_{a<b}(x_b-x_a)$. To describe $\mathsf L$ let
\[
 B(z)=\prod_a\frac{z-x_a}{1-x_az},\qquad \cK_{ab}=\frac1{1-x_ax_b},
\]
be the Blaschke product of the nodes and the Pick (Cauchy) matrix of the Szeg\H{o} kernels $k_a(z)=1/(1-x_az)$.
Put $w_a=1$ at a negative node and $w_a=-5$ at a positive node, and $D=\operatorname{diag}(w_aB'(x_a)x_a^e)$.
Then $\det\mathsf L=\det(\cK+D)/\det\cK$. In the basis $k_a$ the operator $\mathsf L$ acts on the model space
$\cK_B=H^2\ominus BH^2$, which has dimension $n$. The energy method uses only $\|\mathsf L\|\le6$, which gives
$|\det\mathsf L|\le6^n$. On the configurations near the extremal one, the true size of $\det\mathsf L$ is
smaller by a factor $e^{-cN^2}$. We now capture part of this.

Let $k$ and $q$ be the numbers of negative and positive nodes. For a weight $\omega\in\R$ let
$F(\omega)=\cK+\operatorname{diag}(D_-,\omega D_+)$, where $D_-$ and $D_+$ are the diagonal blocks of
$\operatorname{diag}(B'(x_a)x_a^e)$ at the negative and the positive nodes. So $\det\mathsf L=\det F(-5)/\det\cK$. The case $\omega=1$ is a
Pick matrix.

\begin{lemma}\label{lem:pick}
For every configuration $0\le F(1)\le2\cK$.
\end{lemma}

\begin{proof}
For a real vector $c$ put $\phi(z)=B(z)\sum_ac_a/(z-x_a)$. This is the function of $\cK_B$ with node values
$\phi(x_a)=c_aB'(x_a)$, and $\|\phi\|_{H^2}^2=c^T\cK c$. By the residue theorem on the unit disc
\[
 c^T(F(1)-\cK)c=\sum_a\frac{x_a^e\phi(x_a)^2}{B'(x_a)}=\frac1{2\pi i}\oint_{|z|=1}\frac{z^e\phi(z)^2}{B(z)}\,dz .
\]
On the unit circle $|z^e\phi^2/B|=|\phi|^2$, so the right side lies in $[-c^T\cK c,\,c^T\cK c]$.
\end{proof}

\subsection{A Schur complement bound}\label{ssec:schur}
The determinant $\det F(-5)=\det(F(1)-6\operatorname{diag}(0,D_+))$ is a perturbation of a Pick matrix in the
positive block only. The following inequality bounds it by the negative block.

\begin{theorem}\label{thm:C3}
Let $M=\begin{pmatrix}A&C^*\\C&E\end{pmatrix}$ be a Hermitian block matrix with $0\le M\le b$ and $A>0$. Let $T$
be Hermitian of the size of $E$, with $-\tau\le T\le\tau$, and let $c\ge0$. Then
\[
 \Bigl|\det\begin{pmatrix}A&C^*\\C&E-cT\end{pmatrix}\Bigr|\le(b+c\tau)^{\dim E}\det A .
\]
\end{theorem}

\begin{proof}
The determinant equals $\det A\cdot\det(S-cT)$, where $S=E-CA^{-1}C^*$ is the Schur complement. Since $M\ge0$
one has $S\ge0$, and $S\le E\le b$. So $-c\tau\le S-cT\le b+c\tau$, and every eigenvalue of $S-cT$ has modulus at
most $b+c\tau$.
\end{proof}

We apply this in the metric of $\cK$, that is, after the congruence by a Cholesky factor of $\cK$. By
Lemma~\ref{lem:pick} the hypotheses hold with $b=2$ and $c=6$, provided that $\pm D_+\le\tau\cK$ as quadratic
forms. The result is
\begin{equation}\label{eq:C3}
 |\det\mathsf L|\le(2+6\tau)^q\,\det\mathsf L_{--},\qquad
 \det\mathsf L_{--}=\frac{\det(\cK_{--}+D_-)}{\det\cK_{--}},
\end{equation}
where $\cK_{--}$ is the block of $\cK$ at the negative nodes. The factor $(2+6\tau)^q$ is $e^{O(N)}$. All the
information of order $N^2$ is in $\det\mathsf L_{--}$, which depends on the positive nodes only through the
values $B_+(x_b)$ of their Blaschke product at the negative nodes.

\subsection{A contour bound}\label{ssec:T}
It remains to bound $\tau$. The quadratic form of $D_+$ is the sum of the residues of $z^e\phi^2/B$ at the
positive nodes. These lie in the right half of the disc. We integrate over the boundary of the right
half-disc instead of the circle.

\begin{lemma}\label{lem:T}
Let $M_*=\sup_{-1<y<1}|y|^e/|B(iy)|$. For every real vector $c$
\[
 \Bigl|\sum_{x_a>0}B'(x_a)x_a^ec_a^2\Bigr|\le\Bigl(1+\frac{M_*}2\Bigr)c^T\cK c .
\]
If $\frac en-\frac1n\sum_a\Bigl(\frac{t^2}{t^2+x_a^2}-\frac{x_a^2t^2}{1+x_a^2t^2}\Bigr)>0$ for every
$t\in(0,1]$, then $M_*=1$, and the bound is $\frac32c^T\cK c$.
\end{lemma}

\begin{proof}
With $\phi$ as in the proof of Lemma~\ref{lem:pick}, the sum is the contour integral of $z^e\phi^2/B$ over the
boundary of the right half-disc divided by $2\pi i$. No node lies on this boundary. On the arc the integrand
has modulus $|\phi|^2$, and the arc contributes at most $\|\phi\|^2$. On the diameter $z=iy$ the modulus is at
most $M_*|\phi(iy)|^2$. By the Fej\'er--Riesz inequality on a diameter, $\int_{-1}^1|\phi(iy)|^2dy\le\pi
\|\phi\|^2$. So the diameter contributes at most $\frac12M_*\|\phi\|^2$.

For the second statement write $\log(|y|^e/|B(iy)|)=nh(y)$. Then
$h(y)=\frac en\log|y|+\frac1{2n}\sum_a\log\frac{1+x_a^2y^2}{y^2+x_a^2}$, and $h(1)=0$ because $|B(i)|=1$.
Differentiation gives $th'(t)=\frac en-\frac1n\sum_a(\frac{t^2}{t^2+x_a^2}-\frac{x_a^2t^2}{1+x_a^2t^2})$. If
this is positive on $(0,1]$, then $h<0$ on $(0,1)$. Since $h$ is even, $M_*=1$.
\end{proof}

The formal proof avoids the residue theorem. It subtracts the poles with divided differences, integrates
the remaining holomorphic function over the half-disc with primitives, and proves the Fej\'er--Riesz step
by the same device on the upper half-disc. With $\tau=\frac32$, \eqref{eq:C3} becomes
\begin{equation}\label{eq:eleven}
 |\det\mathsf L|\le11^q\det\mathsf L_{--}.
\end{equation}
The condition on $t$ asks that the node measure is not too concentrated near $0$ at any scale. Its left side
at the extremal measure is at least $0.33$ for every $t$.

\subsection{Heine's formula for the negative block}\label{ssec:heine}
For $s\ge1$ and a positive continuous function $f$ on $[-1,1]$ put
\[
 Z_s(f)=\int_{[0,\pi]^s}\prod_{l<l'}(\cos\theta_l-\cos\theta_{l'})^2\prod_lf(\cos\theta_l)\,d\theta .
\]
By Heine's formula $Z_s(f)$ is $s!$ times the Hankel determinant of the moments of $f(u)(1-u^2)^{-1/2}du$. Let
\[
 \phi(u)=\prod_{x_b<0}(1+x_b^2+2x_bu),\qquad \psi(u)=\prod_{x_p>0}(1+x_p^2+2x_pu)^{-1}.
\]

\begin{theorem}\label{thm:heine}
Let $s=q-g$. If $s\ge1$ and $k\le e+1$, then
\[
 \det\mathsf L_{--}=\prod_{b<c}(1-x_bx_c)\prod_{b,p}(1-x_bx_p)^{-1}\cdot\frac{Z_s(\phi\psi)}{Z_s(\psi)},
\]
where $b,c$ run over the negative nodes and $p$ over the positive nodes.
\end{theorem}

\begin{proof}[Sketch]
One first writes $\det\mathsf L_{--}=\det[x_b^j+B_+(x_b)x_b^{m-j}]_{b,j<k}/V(x_-)$ with $m=e+k-1$, by
Lagrange interpolation at the negative nodes. Since $1+x^2+2xu=x\bigl(x+x^{-1}+2u\bigr)$, the factor $\phi$
is a polynomial in $u$ with roots $-\frac12(x_b+x_b^{-1})$. By Christoffel's formula,
$Z_s(\phi\psi)/Z_s(\psi)$ is a determinant of the monic orthogonal polynomials of the weight
$\psi(u)(1-u^2)^{-1/2}$ at these roots, divided by a Vandermonde determinant. That weight is of
Bernstein--Szeg\H{o} type: $\psi$ is the reciprocal of $|h(e^{i\theta})|^2$ for the polynomial
$h(z)=\prod_p(1+x_pz)$. Its orthogonal polynomials of degree $\ge q$ are explicit. At
$u=-\frac12(x+x^{-1})$ they equal Laurent expressions in $x$ built from $h$. Since $h(z)/z^q\bar h(1/z)$
evaluated at $-x_b$ is $B_+(x_b)$ up to sign, the two determinants agree after the Cauchy-type products are
collected.
\end{proof}

The formal proof works with the angle variables throughout. It needs Heine's formula, Andr\'eief's identity
and the Bernstein--Szeg\H{o} polynomials, but no integration over orthogonal groups. Both integrals in
Theorem~\ref{thm:heine} are of log-gas type. There are $s$ particles $u_l=\cos\theta_l$ on $[-1,1]$ with
logarithmic repulsion, in the external field $\log\phi+\log\psi$ or $\log\psi$. Each field is a sum of
logarithmic potentials $l(x,u)=\log(1+x^2+2xu)$ of the nodes.

\subsection{Bounds for the two integrals}\label{ssec:loggas}
The numerator is bounded from above by the method of OpenAI's energy estimate, now with one population. The
key identity is $\log|u-v|=-\log2-\sum_{k\ge1}\frac2kT_k(u)T_k(v)$ for $u\ne v$ in $[-1,1]$. With the moments
$M_k=\sum_lT_k(u_l)$ the pair sum is $-\binom s2\log2-\sum_k\frac1k(M_k^2-\sum_lT_k(u_l)^2)$. For any sequence
$a$ and any $\nu>0$ one has $-M_k^2\le-2\nu a_kM_k+\nu^2a_k^2$. This completes the square against $a$.

\begin{lemma}\label{lem:LGup}
Let $\nu\ge s$ and let $a$ be a summable sequence. With $\mathsf T(a,u)=\sum_ka_kT_k(u)/k$ and
$\|a\|^2=\sum_ka_k^2/k$,
\[
 \log Z_s(f)\le-s^2\log2+2\nu^2\|a\|^2+s\sup_{[-1,1]}\bigl(\log f-4\nu\mathsf T(a,\cdot)\bigr)+o(\nu^2).
\]
\end{lemma}

The diagonal terms $\sum_lT_k(u_l)^2/k$ diverge when summed over $k$. They are regularized by a damping
$\tau^k$ with $\tau=1-e_N$ as in \cite[Prop.~6.1]{OpenAI2026}, at a cost $o(\nu^2)$. The completed squares
$(M_k-\nu a_k)^2/k$ are dropped in this lemma. They are kept in \S\ref{ssec:class}.

The denominator is bounded from below by Jensen's inequality. Let $\nu'$ be a probability measure on $[0,\pi]$
of the form $\nu'=\bigl(1+2\sum_{k\le K}b_k\cos k\theta\bigr)\frac{d\theta}\pi$, with a positive density. In the
variable $u=\cos\theta$ it is a polynomial multiple of the arcsine law.

\begin{lemma}\label{lem:LGlow}
For every bounded measurable $Q$,
\[
 \log Z_s(e^Q)\ge s(s-1)\Bigl(-\log2-\sum_k\frac{2b_k^2}k\Bigr)+s\int Q(\cos\theta)\,d\nu'(\theta)+s\,h(\nu'),
\]
where $h(\nu')$ is the entropy of the density of $\nu'$ with respect to $d\theta$.
\end{lemma}

\begin{proof}
Apply Jensen's inequality to the product density $\prod_l\frac{d\nu'}{d\theta}(\theta_l)$. The logarithmic
energy of $\nu'$ in the variable $u$ is $-\log2-\sum_k2b_k^2/k$, by the expansion of $\log|u-v|$ above.
\end{proof}

The fields of the nodes integrate against $\nu'$ in closed form. For $|x|<1$ one has
$\int l(x,\cos\theta)\,d\nu'=-2\sum_k(-x)^kb_k/k$, a polynomial in $x$. The arcsine part contributes $0$,
because the logarithmic potential of the arcsine law is constant on $[-1,1]$.

\subsection{Positive nodes near $1$}\label{ssec:spikes}
A positive node $x_p=1-\epsilon$ puts into $\psi$ a factor $(1+x_p^2+2x_pu)^{-1}$ of size $\epsilon^{-2}$ near
$u=-1$. Lemma~\ref{lem:LGup} charges the supremum of the field to all $s$ particles. A single such node
would therefore cost about $N\log(1/\epsilon)$, while the true change of $\det\mathsf L_{--}$ is bounded. The
following reduction removes such nodes from the log gas.

\begin{theorem}\label{thm:spike}
Let $\phi$ be positive and non-increasing on $[-1,1]$, and let $\psi=\psi_r\prod_{p\in P}k_p$ with $\psi_r>0$,
$k_p(u)=(1+2y_pu+y_p^2)^{-1}$ and $y_p\in(0,1)$. If $|P|=2m$ or $|P|=2m-1$, then
\[
 \log\frac{Z_s(\phi\psi)}{Z_s(\psi)}\le m\log\phi(-1)+\log\frac{Z_{s-m}(\phi\psi_r)}{Z_{s-m}(\psi_r)} .
\]
\end{theorem}

\begin{proof}[Sketch]
For an odd number of factors add one more factor, at a point to be chosen. Write both integrals as Hankel
determinants. Moving all $y_p$ to a common value changes the ratio monotonically. This uses a one-point
comparison for the weight $(1+2yu+y^2)^{-1}$, whose Chebyshev moments form the geometric sequence $(-y)^k$.
For a common value $y$ the $2m$ factors behave like $m$ particles pinned near $u=-1$. Fischer's inequality
bounds the numerator by $\phi(-1)^m$ times the determinant of the remaining block. As $y\to1$ the Schur
complement of the pinned block in the denominator converges to the Hankel matrix of the remaining
$s-m$ particles.
\end{proof}

We apply the theorem to the positive nodes with $1-x_p<\theta_0=10^{-3}$. Since $\phi(-1)=\prod_b(1-x_b)^2\le
4^k$, each such pair costs at most $O(N)$. The energy of OpenAI's estimate charges each such node
$\beta_hn\log(1/\theta_0)$, which is larger. So the number of such nodes on the configurations that matter
is bounded by a constant times $N$, with a cost absorbed in the margin.

\subsection{Concavity: an inequality for signed logarithmic energy}\label{ssec:reflection}
We now combine OpenAI's energy with \eqref{eq:eleven}, Theorem~\ref{thm:heine} and
Lemmas~\ref{lem:LGup} and~\ref{lem:LGlow}. The total is a functional of the node measure $\mu$, split into its
negative part $\mu_-$ and positive part $\mu_+$, and of the column measure. Its quadratic part is not obviously
negative. Theorem~\ref{thm:heine} contributes the Cauchy-type term $-\iint\log(1-xy)\,d\mu_-(x)d\mu_+(y)$,
and Lemma~\ref{lem:LGup} applied with a trial that depends on $\mu$ contributes a convex term. Write
\[
 I(\delta)=-\iint\log|x-y|\,d\delta(x)d\delta(y),\qquad L(\tau)=-\iint\log(1-xy)\,d\tau(x)d\tau(y)
\]
for real measures on $[-1,1]$ and $[0,1]$. Both are sums of squares of moments: $L(\tau)=\sum_k\frac1k\bigl(\int
x^kd\tau\bigr)^2$. The quadratic part of the combined functional is $-\frac12I-\frac12\sum_k(B_k-2C_k)^2/k+\frac12
L(\mu_+)$, where $B_k$ and $C_k$ are the power moments of the nodes and the columns. The term $\frac12L(\mu_+)$
is absorbed by the following theorem.

\begin{theorem}\label{thm:reflection}
Let $\delta=\delta_-+\delta_+$ be a real measure of finite logarithmic energy and total mass $0$, with
$\delta_-$ on $[-1,0]$ and $\delta_+$ on $[0,1]$. Let $R$ be the reflection $x\mapsto-x$. Then
\[
 I(\delta)\ge L(\delta_+)+L(R\delta_-)\ge L(\delta_+).
\]
\end{theorem}

\begin{proof}
Put $p=\delta_+$ and $r=R\delta_-$, both on $[0,1]$, and let $K_1(a,b)=-\log|a-b|$ and $K_2(a,b)=-\log(a+b)$.
Since $|x-y|=a+b$ for nodes of opposite signs, $I(\delta)=K_1(p)+K_1(r)+2K_2(p,r)$, which equals
$\frac12(K_1+K_2)(p+r)+\frac12(K_1-K_2)(p-r)$.

First, $p+r$ has mass $0$. One has $K_1=G_{\mathbb D}+L$, where $G_{\mathbb D}(a,b)=\log|(1-ab)/(a-b)|$ is the
Green function of the unit disc, whose energy is non-negative. By Frullani's integral
$-\log(a+b)=\int_0^\infty(e^{-(a+b)t}-e^{-t})\frac{dt}t$, so $K_2(\sigma)=\int_0^\infty(\int e^{-at}d\sigma(a))^2
\frac{dt}t\ge0$ for $\sigma$ of mass $0$. Hence $(K_1+K_2)(p+r)\ge L(p+r)$.

Second, $H=K_1-K_2=\log\frac{a+b}{|a-b|}$ satisfies $H(\sigma)\ge L(\sigma)$ for every real $\sigma$ on $[0,1]$.
With $a=e^{-x}$ and $b=e^{-y}$ one has $H=\log\coth\frac{|x-y|}2=\sum_{j\text{ odd}}\frac2je^{-j|x-y|}$. For
$x,y\ge0$, $e^{-j|x-y|}=e^{-j(x+y)}+2j\int_0^{\min(x,y)}e^{-j(x+y-2v)}dv$. Put $F_j(v)=\int_{x\ge v}e^{-j(x-v)}
d\sigma$. Then $H(\sigma)=\sum_{j\text{ odd}}\bigl(\frac2jF_j(0)^2+4\int_0^\infty F_j(v)^2dv\bigr)$. In
$L(\sigma)=\sum_k\frac1k(\int e^{-kx}d\sigma)^2$ the terms with odd $k$ are $F_k(0)^2/k$. For even $k=j+1$,
integration by parts gives $\int e^{-(j+1)x}d\sigma=F_j(0)-\int_0^\infty e^{-(j+1)v}F_j(v)\,dv$. The
Cauchy--Schwarz inequality with weights $1+\frac1j$ and $1+j$ gives $(\int e^{-(j+1)x}d\sigma)^2\le
(1+\frac1j)F_j(0)^2+\frac12\int F_j^2$. Summing, $L(\sigma)\le H(\sigma)$.

Together, $I(\delta)\ge\frac12L(p+r)+\frac12L(p-r)=L(p)+L(r)$.
\end{proof}

By Theorem~\ref{thm:reflection} the quadratic part is at most
$-\frac14I-\frac12\sum_k(B_k-2C_k)^2/k\le0$ on differences of admissible measures. So the combined functional
is concave. Its maximum is bounded by the value at a trial plus the supremum of a linear functional, that
is, by a constant plus the suprema of fields of one variable. The formal proof uses a discrete form of
Theorem~\ref{thm:reflection} for weighted point configurations against a reference measure with bounded
density. The fixed rows of the next subsection change the masses of the two species. The missing mass is
compensated by phantom particles distributed by the arcsine law, whose potential is constant on $[-1,1]$.

\subsection{Fixed rows}\label{ssec:fixedrows}
With $m$ fixed columns the determinant of the row functions is evaluated at the free points and at the fixed
points. The fixed points enter $\det\mathsf L$ as nodes, but they are not integrated over. Lemma~\ref{lem:T}
then needs a change. Its constant $\frac32$ is false when the sum runs over the free positive nodes only, while
the Blaschke product and the Gram form contain the fixed nodes as well. For two adjacent positive nodes $0.5$
and $0.55$ with $e=2$ the best constant is $3.625$. The correct statement keeps the fixed positive residues as a
separate term $W=\sum x_a^e/((1-x_a^2)|B'(x_a)|)$ over the fixed positive nodes. Each of them is bounded by the
reproducing estimate $(c_aB'(x_a))^2(1-x_a^2)\le c^T\cK c$.

Instead of absorbing $W$ into $\tau$, we expand. Write $\det(X+\operatorname{diag}\delta)$ as a sum of principal
minors over subsets $S$ of the fixed positive nodes. This gives
\[
 |\det\mathsf L|\le11^q\sum_S\Bigl(\frac6{11}\Bigr)^{|S|}\det\mathsf L_{--}^{(S)}\,W_S,
\]
where in the term for $S$ the nodes of $S$ interact with the other nodes through $\log(1-xy)$ instead of
$\log|x-y|$. The energy then has three fixed populations: fixed negative nodes, fixed positive nodes, and the
nodes of $S$. Their fields are bounded on their domains like the residual field of Section~\ref{sec:real}.
The fixed nodes are images of quadrature nodes. So the mass of their measure $\kappa$ near $1$ is limited
by the spacing of the quadrature nodes, namely $\kappa([x,1])\le C_1(1-x)$ for $x\ge x_c=0.999$. For such a measure of mass $\alpha$
and a field $\psi$ with $\psi\le\lambda$ on $[0,x_c]$ and $\psi(x)\le g_0+r\log\frac1{1-x^2}$ on $[x_c,1)$,
\[
 \int\psi\,d\kappa\le\alpha\lambda+C_1\int_0^{1-x_c}\bigl(g_0+r\log(1/u)-\lambda\bigr)_+\,du .
\]
The integral has a closed form. This bounds the fixed rows without a finite optimization over their positions.

\subsection{The configurations that matter}\label{ssec:class}
The reduction above needs hypotheses on the configuration. Lemma~\ref{lem:T} needs the condition on $t$. The
number of positive nodes within $\theta_0$ of $1$ must be bounded by a constant times $N$. The fraction of
negative nodes must lie in a fixed window. These hypotheses hold on every configuration whose energy, in
OpenAI's estimate, is within the gain of the maximum. We use the following dichotomy. Either OpenAI's bound at
the rate $U_0$ holds for the configuration, and nothing is to be proved. Or the configuration exceeds it. Then
its energy is within $U_0-U$ of the maximum of OpenAI's bound.

For such a configuration, keep the completed squares and the non-negative differences between the fields and
their suprema in the proof of OpenAI's estimate. Their sum is at most the deficit, a constant times $N^2(U_0-U)$.
This bounds the fraction of nodes on which a field is far below its supremum, by Chebyshev's inequality. It
also bounds every linear statistic $\sum_kF_k(M_k-\nu a_k)/k$ of the moments, by the Cauchy--Schwarz
inequality. The condition of Lemma~\ref{lem:T} is a family of such linear statistics in $t$. It is checked on a
grid of $24$ values of $t$, with a polynomial lower bound between the grid points. The count of nodes near $1$
follows from the field of the positive nodes near $x=1$, which lies below its supremum by a fixed amount. The
window for the fraction of negative nodes follows from polynomial minorants of the indicator functions of the
two signs.

\subsection{The trials and their verification}\label{ssec:signedcert}
The concave majorant of \S\ref{ssec:reflection} is bounded at a trial. A trial consists of a reference measure
for the free nodes, with a piecewise polynomial density, trial sequences for the rows, the columns and the
numerator of Theorem~\ref{thm:heine}, and the measure $\nu'$ of Lemma~\ref{lem:LGlow}. The trial sequences
have $120$ Chebyshev coefficients and geometric tails with complex poles, as in \eqref{eq:trial}. The bound
then reads
\begin{multline*}
 \text{constant}+\alpha\sup\text{(fixed-row fields)}+\sup\text{(free-node fields)}\\
 +\sup\text{(column field)}\le-U_0+\tfrac12\log2+c_2\alpha .
\end{multline*}
The constant is the value at the trial. It consists of logarithmic energies and moments of the reference
measure, the norms of the trial sequences and the entropy of $\nu'$. All of these have closed forms.

The fraction $\alpha$ of fixed columns varies in $[0,\frac1{16}]$. We use trials at $\alpha=0$, $\frac1{200}$ and
$\frac1{16}$. On each of the two intervals between them we interpolate the trials linearly in $\alpha$, in all
the data weighted by mass. Every field is then exactly affine in the interpolation parameter. So its supremum
on the interval is bounded by the chord of its suprema at the two ends. The constant is a quadratic polynomial
in the parameter. Its curvature is an explicit combination of the four cross energies of the two end trials.
The spike constant of Theorem~\ref{thm:spike} interpolates in the same way. At each end trial we therefore
bound the suprema of the following fields: the fields of the negative and the positive free nodes, the
column field, and the fields of the fixed negative, fixed positive and expanded fixed nodes. Each is a
polynomial plus finitely many terms $P(x)\log|x-c|$ and $\operatorname{Re}\rho\log(x-z)$ with complex $z$, plus
a constant.

The suprema are verified on covers by cells, as in \S\ref{ssec:cells}, with the endpoints of the domains
treated by Laurent expansions and logarithmic end windows. The polynomial part of a field has degree about
$120$. On a block of cells we shift its Taylor expansion exactly to order $K$ and bound the remaining
coefficients by one table shared by the block. This keeps the integers small. The three end trials need
\NcellsSigned{} cells in total. The smallest margin of the budget inequality is $0.037$ per unit of $N^2$.

This completes the proof of the bound of the second case at the rate $U_0$ for $\alpha\le\frac1{16}$. With
\S\ref{ssec:cases} it proves Theorem~\ref{thm:real}.
\section{Proof of Theorem~\ref{thm:main}}\label{sec:proof}

\subsection{The criterion}
We prove Proposition~\ref{prop:criterion}. The proof shows where each of the three properties of $Q_N$
enters. The nonvanishing and the smallness are used once, at the end. The divisibility by squares of primes
is what makes the choice of the prime free.

\begin{proof}[Proof of Proposition~\ref{prop:criterion}]
Let $a/b$ be a rational number with $b$ large. Put $L=\log b$, $U=\kappa L$ and $V=(\kappa-\omega)L$, and
suppose $|\xi-a/b|\le b^{-\rho\kappa}$. By the prime number theorem, $(1-\delta)t\le\theta(t)\le(1+\delta)t$
for all large $t$, and $V\to\infty$ with $b$. If every prime $p\in(V,U]$ had $p^2\mid b$, then
$2(\theta(U)-\theta(V))\le\log b=L$. But
\[
 2\bigl(\theta(U)-\theta(V)\bigr)\ge2L\bigl((1-\delta)\kappa-(1+\delta)(\kappa-\omega)\bigr)
 =2L\bigl(\omega-\delta(2\kappa-\omega)\bigr)>L
\]
by the fourth hypothesis. So there are primes in $(V,U]$ whose squares do not divide $b$. Let $p$ be the
\emph{largest} one. Every prime $q\in(p,U]$ has $q^2\mid b$, and $q\le U<2V<2p$ because $2\omega<\kappa$.
Let $B$ be the product of these $q^2$. Then $B\mid b$ and $\log B=2(\theta(U)-\theta(p))$.

Consider the integer $b^{dp}Q_p(a/b)=\sum_ic_ia^ib^{dp-i}$. Each of its terms is divisible by $q^{2dp}$ for
each of these $q$. Indeed $q^{2i}$ divides $c_i$ by hypothesis, and $q^2\mid b$ gives $q^{2(dp-i)}\mid b^{dp-i}$.
So $B^{dp}$ divides the integer. Since $p^2\nmid b$, $Q_p(a/b)\ne0$. So $b^{dp}Q_p(a/b)/B^{dp}$ is a nonzero
integer, and
\begin{equation}\label{eq:recycled}
 1\le\Bigl(\frac bB\Bigr)^{dp}\,|Q_p(a/b)| .
\end{equation}
Since $\rho p\le\rho U=\rho\kappa L$, we have $|a/b-\xi|\le e^{-\rho p}$. So $|Q_p(a/b)|\le e^{-\eta p^2}$, and
\eqref{eq:recycled} gives $\eta p^2\le dp\bigl(L-2(\theta(U)-\theta(p))\bigr)$. With
$\theta(U)-\theta(p)\ge U-p-2\delta U$ this is
\[
 \eta p\le d\bigl(L-2U+2p+4\delta U\bigr).
\]
Hence $(\eta-2d)p\le dL-2dU+4d\delta U$. As $\eta<2d$ and $p\le U$, the left side is at least $(\eta-2d)U$. So
$\eta U\le dL+4d\delta U$, that is, $\eta\kappa\le d+4d\delta\kappa$. This contradicts the last hypothesis.
\end{proof}

The largest admissible prime is the right choice. The primes of the window above it have their squares in
$b$, and the divisibility hypothesis recovers these squares from the coefficients. A cruder choice, of any
prime in a short interval that does not divide $b$, loses a factor. The free lower end $V=(\kappa-\omega)L$
allows $\kappa<2$, because the condition $q\le U<2p$ for the recycled primes needs only $U<2V$.

\subsection{The hypotheses for the design $D$}
Take the design $D=(58,9,3,5,3)$ and $d=53$. The polynomials $Q_p$ of the proposition are the $Q_N$ with
$N=p$. Let
\begin{gather*}
 \rho=\Rho,\quad \eta=\EtaTot,\quad \omega=\tfrac35,\quad \delta=2.9\cdot10^{-7},\\
 \kappa=\frac{410}{\rho}=\Kappa\ldots
\end{gather*}
We check the hypotheses.
\begin{itemize}
\item $\deg Q_N\le53N$ by \eqref{eq:QN} and Lemma~\ref{lem:degree}.
\item The nonvanishing. By Proposition~\ref{prop:nonvanishing}, $\Delta_p(a/b)\ne0$ for large $p$ with
$p^2\nmid b$. Hence $Q_p(a/b)\ne0$.
\item The divisibility. By Lemma~\ref{lem:recycling}, $q^{2i}$ divides the $i$-th coefficient of $Q_p$ for all
primes $p<q\le2p$.
\item The smallness. By Theorem~\ref{thm:real}, Proposition~\ref{prop:denominators} and
$\log(53N)!=o(N^2)$, for $|x-G|\le e^{-\rho N}$ and large $N$,
\[
 |Q_N(x)|\le\exp\bigl(-(\TDelta\ldots-\Aden\ldots-3\varepsilon)N^2\bigr)\le e^{-\eta N^2},
\]
since $\eta=\EtaTot<\EtaExact\ldots$ In the formal proof this is the inequality
$\eta<n_0^2(U_0-\frac12\log2)-\mathsf A$ with $\int\kappa$ replaced by its lower bound $\ChiLean$.
\item The numerical conditions. $\rho\kappa=410$ and $0<\eta<2d=106$. Next $\omega=3/5\le1$ and
$2\omega=1.2<\kappa$. Further $2(\omega-\delta(2\kappa-\omega))=1.199998\ldots>1$. Finally
$d+4d\delta\kappa=\PayL\ldots<\eta\kappa=\PayR\ldots$
\end{itemize}
Proposition~\ref{prop:criterion} therefore gives $|G-a/b|>b^{-\rho\kappa}=b^{-410}$ for all integers $a$ and all
sufficiently large $b$. This proves Theorem~\ref{thm:main}. The last condition is equivalent to
$410>x=d\rho/(\eta-4d\delta)=\Bx\ldots$

\section{Remarks on the method}\label{sec:remarks}

\subsection{The three ingredients}
The exponent $d\rho/\eta$ has three ingredients. The degree $d=n-g$ comes from the construction. Of
$\rho=\Rho$, the energy cost of the fixed columns $n_0c_2$ is $208.62$, and the weight and the binomial factor
of one column contribute $\Psi_0=\PsiZero$. Bounding the two separately, by their maxima $52.79$ and
$58\log2=40.20$, would give $\rho=\RhoOld$. The decay rate $\eta$ is a small difference of two large numbers.
The rank deficiencies of Section~\ref{sec:content} contribute $22.05$ of $\eta=\EtaTot$, and the signed bound
of Section~\ref{sec:signed} contributes $\Gain$. Since the exponent is inversely proportional to $\eta$, a further
gain of $0.1$ in $\eta$ lowers it by about $1.1$.

\subsection{The determinant of the row functions}
The energy method bounds $|\det\mathsf L|$ by $6^n$, which is $e^{O(N)}$ and invisible at the scale $N^2$. At
the extremal configuration, numerical evaluation for small $N$ suggests that $\log|\det\mathsf L|$ is about
$-28N^2$, so most of the true size is not captured. The bound \eqref{eq:eleven} loses only $e^{O(N)}$, and
Theorem~\ref{thm:heine} is an identity. The loss is in the log-gas bounds of \S\ref{ssec:loggas}, which use fixed
trial measures, and in the requirement that the combined functional be bounded by its maximum over all
configurations. A sharper treatment of the denominator $Z_s(\psi)$ would raise the gain. These numerical values
are not part of the formal proof.

\subsection{Designs}
All the results of Section~\ref{sec:content} are lower bounds for the nullity, and the proofs give the kernel
vectors explicitly. The closed formulas, the counting inequality and the families derived from it are proved
for every design. So are the other arguments of the paper, except the stored subspaces of
Section~\ref{ssec:packets} and the trials, which belong to the design $D$. A different design changes $d$,
$\rho$ and $\eta$ together. We make no claim that $D$ is optimal. Except for one pair of families, the function
$\kappa$ uses no relation between different families. At a prime where windows overlap the exact nullity may
be larger than each bound. An upper bound for the nullity is not needed. More independence between families,
or larger stored subspaces, would increase $\eta$ without changing anything else.

\section{Formalization}\label{sec:formal}

All results of this paper are formalized in Lean~4 \cite{Lean4} with Mathlib \cite{Mathlib}, on top of OpenAI's
formalization of the irrationality of $G$ \cite{OpenAI2026}. The development is available at
\begin{center}
\FormalizationLink.
\end{center}
Theorem~\ref{thm:main} is the theorem
\begin{center}
{\small\texttt{CatalanMeasure.GenD.Signed.catalan\_irrationalityMeasure\_410}}
\end{center}
in \lean{CatalanMeasure/GenD/Signed/Data/Assembly/CellsAll.lean}. Its statement uses only Mathlib. Catalan's
constant is written as the series $\sum_j(-1)^j/(2j+1)^2$, and the conclusion is $|G-a/b|>b^{-410}$ for all
integers $a$ and all sufficiently large $b$. The modules contain no \lean{sorry} and no \lean{native\_decide}.
The command \lean{\#print axioms} checks that only the three standard axioms \lean{propext},
\lean{Classical.choice} and \lean{Quot.sound} are used. The repository also contains a configuration
\lean{comparator.json} for \lean{lake comparator} \cite{Comparator}. It checks that a solution file proves a
challenge statement with these axioms only. The prime number theorem is taken from the PNT+ project
\cite{PNTPlus}. The parts of OpenAI's formalization that are used were converted to Lean's module system and
ported to a recent Lean release. The repository's \texttt{README} and \texttt{PORTING.md} describe the changes.

\subsection{Module layout}
The design-generic part of the development is in the directory \lean{GenD} of \lean{CatalanMeasure}. A design is
data. The generic modules prove every statement for an arbitrary design, and the instance files contain
numerals only.
\begin{description}[leftmargin=1.5em,style=nextline]\sloppy
\item[Generic layer] \lean{PL} (piecewise linear profiles and their areas), \lean{PrimeGrowth} (the weighted
prime number theorem \eqref{eq:weighted-pnt}), \lean{Budget} (the scalar budgets), \lean{PayW} and
\lean{Recycled} (Proposition~\ref{prop:criterion}), and \lean{Final} (the structure \lean{DesignData} and the
theorem \lean{measure\_of\_design}, which contain no numeral).
\item[Denominators] \lean{OddProfile}, \lean{TwoAdic}, \lean{DenomExp}, \lean{DenomGrowth},
\lean{Denominator} and \lean{Normalizer}, with \lean{Uniform/Arithmetic} and \lean{Uniform/LedgerW} (the
common denominator built from an arbitrary lower bound for the nullity, and the growth of its logarithm).
\item[Rank deficiencies] \lean{RankT}, \lean{RankU}, \lean{RankU6}, \lean{RankS} and their shared layers for
Theorem~\ref{thm:content}; \lean{Uniform/W/} for the counting inequality (\lean{Core}, \lean{Profile}),
Lucas's step (\lean{Lucas/PStep}), the families of Section~\ref{ssec:families} and the stored data of
Section~\ref{ssec:packets} (\lean{Prof/Packet}, \lean{Charts6/Check}); the profile of $D$ with its area
\lean{gain532} and the bound \lean{gain532\_ge} in \lean{Uniform/Bank532}.
\item[Nonvanishing] \lean{Nonvanishing}, with the residue-table certificates modulo $1009$.
\item[Real place, perturbation] \lean{Analytic/} (the energy method with fixed points, the column expansion
and the quadrature nodes); \lean{BinomVandermonde} (Lemma~\ref{lem:binomvand}) and \lean{Keyhole}
(Lemma~\ref{lem:keyhole}) in \lean{GenD/Signed}; \lean{D58/Budget590} (the budgets at the rates $U$ and $U_0$,
and $\rho=$ \lean{rhoEnd}); \lean{Signed/Data/S1} (the $14$ cells of the bound for $\Psi_0$).
\item[Real place, signed bound] \lean{GenD/Signed}: \lean{PickDet} and \lean{RealPlace}
(Lemma~\ref{lem:pick}, Theorem~\ref{thm:C3}, \eqref{eq:eleven}), \lean{LemmaT}, \lean{LemmaTFree} and
\lean{HalfDisc} (Lemma~\ref{lem:T}), \lean{RatioIdentity} and \lean{BernsteinSzego}
(Theorem~\ref{thm:heine}), \lean{TiltedTangent} and \lean{ArcsineDenominator} (Lemmas~\ref{lem:LGup}
and~\ref{lem:LGlow}), \lean{Spike} and its companions (Theorem~\ref{thm:spike}), \lean{ReflectionEnergy}
and \lean{ArcsinePhantom} (Theorem~\ref{thm:reflection}), \lean{MarkedExpansion}, \lean{SigmaCap} and
\lean{EnergyNear} (\S\ref{ssec:fixedrows}), \lean{LemmaK} and \lean{Brackets} (\S\ref{ssec:class}),
\lean{Interp} and \lean{Phase2Contract} (the interpolation of trials), and \lean{Gate} and \lean{Theorem410}
(the assembly).
\item[Trials and cells] \lean{Signed/Data/Trial} (the three trials), \lean{Signed/Data/Consts} (the constants
of the trials, evaluated in exact arithmetic), \lean{Signed/Data/Cells} (the cell engine and the cells of the
three trials) and \lean{Signed/Data/Assembly} (the records of the cell statements and the final theorem).
\end{description}

\subsection{Certificates}
The proof combines conventional formal proofs with numerical certificates checked by the Lean kernel. They
are the inverses of the base matrices modulo $1009$, the trial sequences and covers of
Sections~\ref{sec:real} and~\ref{sec:signed}, the cells of the bound for $\Psi_0$, and the stored subspaces,
kernel vectors and window conditions of Section~\ref{ssec:packets}. All are exact rational or integer
computations, decided by the kernel without compiled code. The interval arithmetic used to generate the trial
data and the covers \cite{Arb} is not trusted. Lean rechecks every inequality. Five devices keep the
certificates small enough for the kernel.
\begin{itemize}
\item \emph{Shared Taylor bounds.} The cells of a cover are grouped in blocks. For each block the derivatives of
the polynomial part of the field are evaluated once, exactly, at the centre of the block, together with one
bound for the remainder. Every cell of the block obtains its expansion by an exact shift.
\item \emph{Split shifts.} In the covers of Section~\ref{sec:signed}, only the low-order coefficients of the
shifted expansion are computed exactly. The others are bounded by one crude table per block.
\item \emph{Accuracy per term.} The logarithms are evaluated by the series
$\log\frac{1+x}{1-x}=2x\sum_ix^{2i}/(2i+1)$ with an explicit geometric tail. The number of terms is chosen for
each logarithmic term separately, from a table of nodes shared by all the cells.
\item \emph{Residue recurrences.} The tables of moments modulo $1009$ are generated by the one-step recurrences
that define them, once proved correct. The kernel checks the recurrence and the products of the matrices with
their claimed inverses, row by row, in natural numbers.
\item \emph{$2$-adic normalisation.} Where a recurrence or a shift produces integers divisible by large powers
of $2$, the tables store the quotients by these powers, with one lemma relating the two. This also avoids
integers that the kernel's hashing of literals treats as equal.
\end{itemize}

\section*{Disclosure of AI use}

The mathematics, the Lean formalization, the certificate generators and the text of this paper were produced by
AI systems at the request of the author, who posed the problems, directed the work and is responsible for this
paper: Claude Opus~5.5, Claude Sonnet~5.5 and Claude Haiku~5.5 (Anthropic), running in Claude Code, and
GPT-6.1 Sol (OpenAI), running in Codex. The underlying irrationality proof for $G$ and its Lean formalization are
OpenAI's \cite{OpenAI2026}.

\subsection*{Acknowledgements}
We thank the authors of Mathlib and of the PNT+ project, and OpenAI for making its formalization available.

\providecommand{\bysame}{\leavevmode\hbox to3em{\hrulefill}\thinspace}
\providecommand{\MR}{\relax\ifhmode\unskip\space\fi MR }
% \MRhref is called by the amsart/book/proc definition of \MR.
\providecommand{\MRhref}[2]{%
  \href{http://www.ams.org/mathscinet-getitem?mr=#1}{#2}
}
\providecommand{\href}[2]{#2}
\begin{thebibliography}{10}

\bibitem{Apery79}
Roger Ap\'ery, \emph{Irrationalit\'e de {$\zeta(2)$} et {$\zeta(3)$}},
  Ast\'erisque \textbf{61} (1979), 11--13.

\bibitem{Beukers79}
Frits Beukers, \emph{A note on the irrationality of {$\zeta(2)$} and
  {$\zeta(3)$}}, Bull. London Math. Soc. \textbf{11} (1979), no.~3, 268--272.

\bibitem{Bojnal2026}
Hrishik~Sai Bojnal, \emph{An irrationality exponent below 2402 for {C}atalan's
  constant}, GitHub repository, October 15, 2026,
  \url{https://github.com/proof-factory/catalan-exponent}; design search and
  certification by Claude Fable~5.1 (Anthropic), 2026.

\bibitem{CDT2024}
Frank Calegari, Vesselin Dimitrov, and Yunqing Tang, \emph{The linear
  independence of {$1$}, {$\zeta(2)$}, and {$L(2,\chi_{-3})$}}, 2024,
  \href{https://arxiv.org/abs/2408.15403}{arXiv:2408.15403}.

\bibitem{CDT2025}
\bysame, \emph{Arithmetic holonomy bounds and effective {D}iophantine
  approximation}, 2025, Theorem 5.1 gives $\mu(L(2,\chi_{-3}))\le 24781$,
  \href{https://arxiv.org/abs/2510.04156}{arXiv:2510.04156}.

\bibitem{Carreira2026b}
Marcos Costa Santos~Carreira, \emph{Determinant constructions for the
  {D}irichlet beta value at four}, Preprint, 9 October 2026,
  \href{https://doi.org/10.5281/zenodo.23326647}{doi:10.5281/zenodo.23326647},
  2026.

\bibitem{Carreira2026}
\bysame, \emph{General designs and an irrationality exponent for {C}atalan's
  constant}, Preprint, 9 October 2026, revised 21 October 2026,
  \href{https://doi.org/10.5281/zenodo.23532239}{doi:10.5281/zenodo.23532239};
  sources at \url{https://github.com/MarcosCarreira/catalan-designs}, 2026.

\bibitem{Lean4}
Leonardo de~Moura and Sebastian Ullrich, \emph{The {L}ean 4 theorem prover and
  programming language}, Automated Deduction -- CADE 28, Lecture Notes in
  Computer Science, vol. 12699, Springer, 2021, pp.~625--635.

\bibitem{Eskandari2026}
Payman Eskandari, \emph{Improved rational approximations to {C}atalan's
  constant}, 2026, Preprint, 22 September 2026,
  \href{https://arxiv.org/abs/2609.26354}{arXiv:2609.26354}.

\bibitem{EMN2025}
Payman Eskandari, Kumar Murty, and Yusuke Nemoto, \emph{Mixed motives and
  linear forms in the {C}atalan constant}, 2025, Preprint, October 2025 (v3: 29
  August 2026), \href{https://arxiv.org/abs/2510.20648}{arXiv:2510.20648}.

\bibitem{Fauzan2026}
Aabir Fauzan, \emph{{$\zeta(5)$} is irrational}, Preprint v1, September 17,
  2026, Zenodo,
  \href{https://doi.org/10.5281/zenodo.22826419}{doi:10.5281/zenodo.22826419},
  2026.

\bibitem{Forrester2018Andreief}
Peter~J. Forrester, \emph{Meet {A}ndr\'eief, {B}ordeaux 1886, and {A}ndreev,
  {K}harkov 1882--83}, 2018,
  \href{https://arxiv.org/abs/1806.10411}{arXiv:1806.10411}.

\bibitem{Giraudin26}
David~Niedbala Giraudin, \emph{A new upper bound for the irrationality exponent
  of {$\zeta(2)$}}, 2026, Preprint, 2 October 2026,
  \href{https://arxiv.org/abs/2610.02912}{arXiv:2610.02912}.

\bibitem{Giraudin26z3}
\bysame, \emph{A new upper bound for the irrationality exponent of
  {$\zeta(3)$}}, 2026, Preprint, 22 September 2026,
  \href{https://arxiv.org/abs/2609.26980}{arXiv:2609.26980}.

\bibitem{Hata93}
Masayoshi Hata, \emph{Rational approximations to {$\pi$} and some other
  numbers}, Acta Arith. \textbf{63} (1993), no.~4, 335--349.

\bibitem{Jha2026}
R.~K. Jha, \emph{A finite irrationality measure for {C}atalan's constant},
  Preprint v1, October 8, 2026,
  \href{https://doi.org/10.5281/zenodo.23237269}{doi:10.5281/zenodo.23237269};
  Lean sources at \url{https://github.com/rohitkrjha/catalan-measure}, 2026.

\bibitem{Arb}
Fredrik Johansson, \emph{Arb: efficient arbitrary-precision midpoint-radius
  interval arithmetic}, IEEE Trans. Comput. \textbf{66} (2017), no.~8,
  1281--1292, \href{https://arxiv.org/abs/1611.02831}{arXiv:1611.02831}.

\bibitem{KrattRivoal08}
Christian Krattenthaler and Tanguy Rivoal, \emph{On a linear form for
  {C}atalan's constant}, South East Asian J. Math. Math. Sci. \textbf{6}
  (2008), no.~2, 3--15,
  \href{https://arxiv.org/abs/0810.1927}{arXiv:0810.1927}.

\bibitem{KrattZudilin19}
Christian Krattenthaler and Wadim Zudilin, \emph{Hypergeometry inspired by
  irrationality questions}, Kyushu J. Math. \textbf{73} (2019), no.~1,
  189--203, \href{https://arxiv.org/abs/1802.08856}{arXiv:1802.08856}.

\bibitem{Lai2026}
Li~Lai, \emph{The linear independence of {$1$}, {$\zeta(3)$}, and
  {$\zeta(4)$}}, 2026, Preprint, 5 October 2026 (v2: 15 October 2026),
  \href{https://arxiv.org/abs/2610.06838}{arXiv:2610.06838}.

\bibitem{Comparator}
{Lean FRO}, \emph{Comparator: checking a {L}ean proof against a trusted
  statement}, \url{https://github.com/leanprover/comparator}, 2026, Software;
  accessed October 2026.

\bibitem{LBKWSK26}
Elyasheev Leibtag, Hila Barkan, Rotem Kalisch, Shachar Weinbaum, Michael
  Shalyt, and Ido Kaminer, \emph{New directions in irrationality}, 2026,
  Preprint, 20 October 2026,
  \href{https://arxiv.org/abs/2610.27799}{arXiv:2610.27799}; supplement and
  code at
  \url{https://github.com/RamanujanMachine/new-directions-in-irrationality}.

\bibitem{Litsey2026b}
Sean Litsey, \emph{Irrationality measures for differences of trigamma values},
  Hexagon preprint hexagon:2610.00401,
  \url{https://hexagonmath.org/2610.00401}, 2026.

\bibitem{LiuJiangZhang26}
Jingwen Liu, Kai Jiang, and Pingwen Zhang, \emph{Irrationality exponents of
  logarithms of positive rational numbers}, 2026, Preprint, 7 October 2026,
  revised 8 October 2026 (v2),
  \href{https://arxiv.org/abs/2610.10192}{arXiv:2610.10192}.

\bibitem{Marcovecchio09}
Raffaele Marcovecchio, \emph{The {R}hin--{V}iola method for {$\log 2$}}, Acta
  Arith. \textbf{139} (2009), no.~2, 147--184.

\bibitem{OpenAI2026}
{OpenAI}, \emph{Catalan's constant is irrational}, Preprint dated September 24,
  2026, released October 6, 2026, with a Lean formalization, available at
  \href{https://github.com/openai/math/blob/fd4aeeb2ee4fc729c18d98444fed42fd0529eeeb/preprints/Catalans-constant-is-irrational-September-24-2026/paper.pdf}{the
  original paper PDF} and
  \href{https://github.com/openai/math/tree/fd4aeeb2ee4fc729c18d98444fed42fd0529eeeb/lean/OAI/NumberTheory/Catalan}{the
  Lean formalization} (repository commit \texttt{fd4aeeb}, October 7, 2026),
  2026.

\bibitem{OpenAIPi2026}
\bysame, \emph{The irrationality exponent of {$\pi$} is 2}, Preprint dated
  September 24, 2026, released October 6, 2026,
  \url{https://github.com/openai/math/tree/main/preprints/The-irrationality-exponent-of-pi-is-2-September-24-2026},
  2026.

\bibitem{Prellberg26}
Thomas Prellberg, \emph{Three irrationality results for the dilogarithm:
  {R}hin--{V}iola and {V}iola--{Z}udilin constructions at {$-1/4$}, {$1/5$},
  and {$-1/3$}}, 2026, Preprint, 7 September 2026,
  \href{https://arxiv.org/abs/2609.06895}{arXiv:2609.06895}.

\bibitem{RhinViola96}
Georges Rhin and Carlo Viola, \emph{On a permutation group related to
  {$\zeta(2)$}}, Acta Arith. \textbf{77} (1996), no.~1, 23--56.

\bibitem{RhinViola01}
\bysame, \emph{The group structure for {$\zeta(3)$}}, Acta Arith. \textbf{97}
  (2001), no.~3, 269--293.

\bibitem{RivoalZudilin2003}
Tanguy Rivoal and Wadim Zudilin, \emph{Diophantine properties of numbers
  related to {C}atalan's constant}, Math. Ann. \textbf{326} (2003), no.~4,
  705--721.

\bibitem{SaffTotik}
Edward~B. Saff and Vilmos Totik, \emph{Logarithmic potentials with external
  fields}, Grundlehren der mathematischen Wissenschaften, vol. 316, Springer,
  Berlin, 1997.

\bibitem{Salikhov08}
V.~Kh. Salikhov, \emph{On the irrationality measure of the number {$\pi$}},
  Russian Math. Surveys \textbf{63} (2008), no.~3, 570--572.

\bibitem{Sun2026}
Zhi-Wei Sun, \emph{Catalan's constant is irrational}, 2026,
  \href{https://arxiv.org/abs/2609.04176}{arXiv:2609.04176}.

\bibitem{Mathlib}
{The mathlib Community}, \emph{The {L}ean mathematical library}, Proceedings of
  the 9th ACM SIGPLAN International Conference on Certified Programs and Proofs
  (CPP 2020), 2020, pp.~367--381.

\bibitem{PNTPlus}
{The PrimeNumberTheoremAnd contributors}, \emph{The {PNT+} project: a {L}ean
  formalization of the prime number theorem and related results},
  \url{https://github.com/AlexKontorovich/PrimeNumberTheoremAnd}, The vendored
  compatibility additions used here are recorded in the repository's
  \texttt{NOTICE} and \texttt{PORTING.md}.

\bibitem{Wachs2026}
Dane Wachs, \emph{A note on a recent claimed proof of the irrationality of
  {C}atalan's constant}, 2026,
  \href{https://arxiv.org/abs/2609.22339}{arXiv:2609.22339}.

\bibitem{Xu2026}
Ryan Xu, \emph{Improved {$\kappa=2$} trial: {$\mu(G)\le727.03$}}, GitHub issue,
  October 22, 2026,
  \url{https://github.com/slitsey/catalan-irrationality-measure/issues/5}; data
  and scripts at \url{https://github.com/ry-xu/slightlysupervised}, 2026.

\bibitem{ZeilbergerZudilin20}
Doron Zeilberger and Wadim Zudilin, \emph{The irrationality measure of {$\pi$}
  is at most {$7.103205334137\ldots$}}, Mosc. J. Comb. Number Theory \textbf{9}
  (2020), no.~4, 407--419,
  \href{https://arxiv.org/abs/1912.06345}{arXiv:1912.06345}.

\bibitem{Zudilin02Cheb}
Wadim Zudilin, \emph{A few remarks on linear forms involving {C}atalan's
  constant}, Chebyshevskii Sb. \textbf{3} (2002), no.~2(4), 60--70, In Russian,
  \href{https://arxiv.org/abs/math/0210423}{arXiv:math/0210423}.

\bibitem{Zudilin2003}
\bysame, \emph{An {A}p\'ery-like difference equation for {C}atalan's constant},
  Electron. J. Combin. \textbf{10} (2003), Research Paper 14,
  \href{https://arxiv.org/abs/math/0201024}{arXiv:math/0201024}.

\bibitem{Zudilin14}
\bysame, \emph{Two hypergeometric tales and a new irrationality measure of
  {$\zeta(2)$}}, Ann. Math. Qu\'ebec \textbf{38} (2014), no.~1, 101--117,
  \href{https://arxiv.org/abs/1310.1526}{arXiv:1310.1526}.

\bibitem{Zudilin}
\bysame, \emph{A determinantal approach to irrationality}, Constr. Approx.
  \textbf{45} (2017), no.~2, 301--310,
  \href{https://arxiv.org/abs/1507.05697}{arXiv:1507.05697}.

\bibitem{Zudilin19Survey}
\bysame, \emph{Arithmetic of {C}atalan's constant and its relatives}, Abh.
  Math. Semin. Univ. Hambg. \textbf{89} (2019), no.~1, 45--53,
  \href{https://arxiv.org/abs/1804.09922}{arXiv:1804.09922}.

\end{thebibliography}

\end{document}
